\documentclass[12pt, ukrainian]{article}
\usepackage[a4paper]{geometry}
\setlength\parindent{0mm}
\geometry{top=1.1cm,bottom=1.1cm,left=1.0cm,right=1.0cm}
\pagestyle{empty}
\usepackage[T2A]{fontenc}
\usepackage[utf8]{inputenc}
\usepackage[ukrainian]{babel}
\usepackage{graphicx}
\usepackage{caption}
\DeclareCaptionFormat{nocaption}{}
\captionsetup{format=nocaption,aboveskip=0pt,belowskip=0pt}
\usepackage[Export]{adjustbox}
\adjustboxset{max size={0.9\linewidth}{0.9\paperheight}}
\begin{document}
{\Large{06.11.2025 Контрольна робота \No 1, варіант  \No \textbf { 1 }\par}}
{
\begin {large}
У завданнях 1-9 вказати, що буде виведено за виконання.

Повідомлення інтерпретатора про помилку позначати ERROR

\begin{minipage}[t]{7.5cm}

\begin{center}
\textbf{Завдання 1}
\end{center}

\begin{verbatim}
print(5/7/4*4)
\end{verbatim}



\begin{center}
\textbf{Завдання 2}
\end{center}

\begin{verbatim}
print(1**-0.5*False)
\end{verbatim}



\begin{center}
\textbf{Завдання 4}
\end{center}

\begin{verbatim}
print("1" >= "9.0j")
\end{verbatim}



\begin{center}
\textbf{Завдання 6}
\end{center}

\begin{verbatim}

def f(n):
    print("f", end="");
    return n>=-1

print(f(-7) or f(7))
\end{verbatim}



\begin{center}
\textbf{Завдання 7}
\end{center}

\begin{verbatim}
def f(a):
    w=85
    if a: 
        return 1
    if a!=0:
         w=0
    else:
         return 5
    return w

print(f(0))
\end{verbatim}

\end{minipage}
\vline \hspace{6mm}
\begin{minipage}[t]{10.5cm}

\begin{center}
\textbf{Завдання 3}
\end{center}

\begin{verbatim}
print(not 2.0<=2 or True!=8.0 or 8>7)
\end{verbatim}



\begin{center}
\textbf{Завдання 5}
\end{center}

\begin{verbatim}
print(6 is True == 7 < 7.0)
\end{verbatim}



\begin{center}
\textbf{Завдання 8}
\end{center}

\begin{verbatim}
a, b, c = 6, 8, 7
def f(a, b, c):
    print(a, b, c, end=" ")

f(4, 2)
print(a, b, c)
\end{verbatim}



\begin{center}
\textbf{Завдання 9}
\end{center}

\begin{verbatim}
a,b,c=9,6,1
def g(a):
    a=3
    b+=1
    c=2
    return a+b+c

a,b,c=7,0,2
print(g(b),a,b,c)
\end{verbatim}



\begin{center}
\textbf{Завдання 10}
\end{center}

Написати функцiю f, що отримує 2 числа і повертає логічну ознаку того, що вони рiвнi з точнiстю до 0.01. Стиль запису умови також оцінюється.\par

\end{minipage}

\end {large}}
\newpage

{\Large{06.11.2025 Контрольна робота \No 1, варіант  \No \textbf { 2 }\par}}
{
\begin {large}
У завданнях 1-9 вказати, що буде виведено за виконання.

Повідомлення інтерпретатора про помилку позначати ERROR

\begin{minipage}[t]{7.5cm}

\begin{center}
\textbf{Завдання 1}
\end{center}

\begin{verbatim}
print(4//3//3*8)
\end{verbatim}



\begin{center}
\textbf{Завдання 2}
\end{center}

\begin{verbatim}
print(4**-1.0*2)
\end{verbatim}



\begin{center}
\textbf{Завдання 4}
\end{center}

\begin{verbatim}
print(4 == 7.0)
\end{verbatim}



\begin{center}
\textbf{Завдання 6}
\end{center}

\begin{verbatim}

def f(n):
    print("f", end="");
    return n

print(f(3) and f(-8))
\end{verbatim}



\begin{center}
\textbf{Завдання 7}
\end{center}

\begin{verbatim}
def h(c):
    x=21
    if c<=1: 
        x=2
    elif c<-1:
         x=4
    else:
         return 9
    return x

print(h(3))
\end{verbatim}

\end{minipage}
\vline \hspace{6mm}
\begin{minipage}[t]{10.5cm}

\begin{center}
\textbf{Завдання 3}
\end{center}

\begin{verbatim}
print(not False==4 and 6>=7.0 and 3<9.0)
\end{verbatim}



\begin{center}
\textbf{Завдання 5}
\end{center}

\begin{verbatim}
print(9.0 >= 9 > False != 3)
\end{verbatim}



\begin{center}
\textbf{Завдання 8}
\end{center}

\begin{verbatim}
a, b, c = 9, 7, 6
def h(a, b=8, c=9):
    print(a, b, c, end=" ")

h(0, c=1, b=3)
print(a, b, c)
\end{verbatim}



\begin{center}
\textbf{Завдання 9}
\end{center}

\begin{verbatim}
a,b,c=8,4,5
def f(b):
    a=4
    b-=5
    c=1
    return a+b+c

a,b,c=3,7,9
print(f(b),a,b,c)
\end{verbatim}



\begin{center}
\textbf{Завдання 10}
\end{center}

Написати функцiю f, що отримує 2 числа і повертає логічну ознаку того, що хоча б одне з них не нуль. Стиль запису умови також оцінюється.\par

\end{minipage}

\end {large}}
\newpage

{\Large{06.11.2025 Контрольна робота \No 1, варіант  \No \textbf { 3 }\par}}
{
\begin {large}
У завданнях 1-9 вказати, що буде виведено за виконання.

Повідомлення інтерпретатора про помилку позначати ERROR

\begin{minipage}[t]{7.5cm}

\begin{center}
\textbf{Завдання 1}
\end{center}

\begin{verbatim}
print(6/8%6*5)
\end{verbatim}



\begin{center}
\textbf{Завдання 2}
\end{center}

\begin{verbatim}
print(2**-2+-2)
\end{verbatim}



\begin{center}
\textbf{Завдання 4}
\end{center}

\begin{verbatim}
print(True < "False")
\end{verbatim}



\begin{center}
\textbf{Завдання 6}
\end{center}

\begin{verbatim}

def f(n):
    print("f", end="");
    return n==-3

print(f(-2) or f(0))
\end{verbatim}



\begin{center}
\textbf{Завдання 7}
\end{center}

\begin{verbatim}
def g(b):
    v=28
    if b==4: 
        v=7
    if b>=3:
         return 3
    else:
         v=8
    return v

print(g(5))
\end{verbatim}

\end{minipage}
\vline \hspace{6mm}
\begin{minipage}[t]{10.5cm}

\begin{center}
\textbf{Завдання 3}
\end{center}

\begin{verbatim}
print(6.0<9 or not 5!=5 and True==3.0)
\end{verbatim}



\begin{center}
\textbf{Завдання 5}
\end{center}

\begin{verbatim}
print(8.0 <= 5 <= 4 is 3.0)
\end{verbatim}



\begin{center}
\textbf{Завдання 8}
\end{center}

\begin{verbatim}
a, b, c = 9, 8, 6
def g(a, b=7, c):
    print(a, b, c, end=" ")

g(a=5, 1, c=5)
print(a, b, c)
\end{verbatim}



\begin{center}
\textbf{Завдання 9}
\end{center}

\begin{verbatim}
a,b,c=6,8,0
def h(b):
    a=3
    b*=4
    c=5
    return a+b+c

a,b,c=3,1,5
print(h(b),a,b,c)
\end{verbatim}



\begin{center}
\textbf{Завдання 10}
\end{center}

Написати функцiю f, що отримує 2 числа і повертає логічну ознаку того, що перше бiльше. Стиль запису умови також оцінюється.\par

\end{minipage}

\end {large}}
\newpage

{\Large{06.11.2025 Контрольна робота \No 1, варіант  \No \textbf { 4 }\par}}
{
\begin {large}
У завданнях 1-9 вказати, що буде виведено за виконання.

Повідомлення інтерпретатора про помилку позначати ERROR

\begin{minipage}[t]{7.5cm}

\begin{center}
\textbf{Завдання 1}
\end{center}

\begin{verbatim}
print(9//4*5*6)
\end{verbatim}



\begin{center}
\textbf{Завдання 2}
\end{center}

\begin{verbatim}
print(2**1.0--1)
\end{verbatim}



\begin{center}
\textbf{Завдання 4}
\end{center}

\begin{verbatim}
print("True" != 6)
\end{verbatim}



\begin{center}
\textbf{Завдання 6}
\end{center}

\begin{verbatim}

def f(n):
    print("f", end="");
    return n

print(f(5) and f(-4))
\end{verbatim}



\begin{center}
\textbf{Завдання 7}
\end{center}

\begin{verbatim}
def h(d):
    z=24
    if d>-4: 
        return 6
    elif d<=-3:
         return 3
    else:
         z=0
    return z

print(h(-3))
\end{verbatim}

\end{minipage}
\vline \hspace{6mm}
\begin{minipage}[t]{10.5cm}

\begin{center}
\textbf{Завдання 3}
\end{center}

\begin{verbatim}
print(False<=5.0 and not 4>=6 or 4.0>2)
\end{verbatim}



\begin{center}
\textbf{Завдання 5}
\end{center}

\begin{verbatim}
print(2.0 < 4.0 == 2 >= True)
\end{verbatim}



\begin{center}
\textbf{Завдання 8}
\end{center}

\begin{verbatim}
a, b, c = 6, 9, 8
def g(a, b, c=7):
    print(a, b, c, end=" ")

g(b=2, c=0, 4)
print(a, b, c)
\end{verbatim}



\begin{center}
\textbf{Завдання 9}
\end{center}

\begin{verbatim}
a,b,c=3,6,4
def h(a):
    a=4
    b=1
    c=5
    return a+b+c

a,b,c=2,1,9
print(h(a),a,b,c)
\end{verbatim}



\begin{center}
\textbf{Завдання 10}
\end{center}

Написати функцiю f, що отримує 2 цілих числа i повертає логічну ознаку того, що їх трійковий запис закінчується на одну ту саму цифру. Стиль запису умови також оцінюється.\par

\end{minipage}

\end {large}}
\newpage

{\Large{06.11.2025 Контрольна робота \No 1, варіант  \No \textbf { 5 }\par}}
{
\begin {large}
У завданнях 1-9 вказати, що буде виведено за виконання.

Повідомлення інтерпретатора про помилку позначати ERROR

\begin{minipage}[t]{7.5cm}

\begin{center}
\textbf{Завдання 1}
\end{center}

\begin{verbatim}
print(7//6/8*7)
\end{verbatim}



\begin{center}
\textbf{Завдання 2}
\end{center}

\begin{verbatim}
print(9**1.0*True)
\end{verbatim}



\begin{center}
\textbf{Завдання 4}
\end{center}

\begin{verbatim}
print("1.0" > False)
\end{verbatim}



\begin{center}
\textbf{Завдання 6}
\end{center}

\begin{verbatim}

def f(n):
    print("f", end="");
    return n<=2

print(f(2) and f(-1))
\end{verbatim}



\begin{center}
\textbf{Завдання 7}
\end{center}

\begin{verbatim}
def h(a,b):
    c=89
    if b==-1:
        c=1
    if b!=0:
         return 7
    else: 
        c=2
    return c

print(h(4,-4))
\end{verbatim}

\end{minipage}
\vline \hspace{6mm}
\begin{minipage}[t]{10.5cm}

\begin{center}
\textbf{Завдання 3}
\end{center}

\begin{verbatim}
print(not 2.0<8 or 9>False and 7.0==3)
\end{verbatim}



\begin{center}
\textbf{Завдання 5}
\end{center}

\begin{verbatim}
print(8 != 7 > 2 <= 5.0)
\end{verbatim}



\begin{center}
\textbf{Завдання 8}
\end{center}

\begin{verbatim}
a, b, c = 8, 9, 6
def f(a, b=7, c):
    print(a, b, c, end=" ")

f(3, 0, 3)
print(a, b, c)
\end{verbatim}



\begin{center}
\textbf{Завдання 9}
\end{center}

\begin{verbatim}
a,b,c=4,2,8
def f(a):
    a=2
    b*=3
    c=1
    return a+b+c

a,b,c=1,3,6
print(f(a),a,b,c)
\end{verbatim}



\begin{center}
\textbf{Завдання 10}
\end{center}

Написати функцiю f, що отримує рядок i повертає логічну ознаку того, що вiн складається рівно з однієї літери, яка є двійковою цифрою. Стиль запису умови також оцінюється.\par

\end{minipage}

\end {large}}
\newpage

{\Large{06.11.2025 Контрольна робота \No 1, варіант  \No \textbf { 6 }\par}}
{
\begin {large}
У завданнях 1-9 вказати, що буде виведено за виконання.

Повідомлення інтерпретатора про помилку позначати ERROR

\begin{minipage}[t]{7.5cm}

\begin{center}
\textbf{Завдання 1}
\end{center}

\begin{verbatim}
print(3/5//9*9)
\end{verbatim}



\begin{center}
\textbf{Завдання 2}
\end{center}

\begin{verbatim}
print(1**0.5*1)
\end{verbatim}



\begin{center}
\textbf{Завдання 4}
\end{center}

\begin{verbatim}
print("5j" <= 2.0)
\end{verbatim}



\begin{center}
\textbf{Завдання 6}
\end{center}

\begin{verbatim}

def f(n):
    print("f", end="");
    return n

print(f(1) or f(6))
\end{verbatim}



\begin{center}
\textbf{Завдання 7}
\end{center}

\begin{verbatim}
def g(b):
    y=70
    if b: 
        y=5
    elif b>=2:
         return 4
    else:
         return 2
    return y

print(g(-8))
\end{verbatim}

\end{minipage}
\vline \hspace{6mm}
\begin{minipage}[t]{10.5cm}

\begin{center}
\textbf{Завдання 3}
\end{center}

\begin{verbatim}
print(not True!=4.0 and 7<=6.0 or 2>=9)
\end{verbatim}



\begin{center}
\textbf{Завдання 5}
\end{center}

\begin{verbatim}
print(4 != 6.0 > False == 6)
\end{verbatim}



\begin{center}
\textbf{Завдання 8}
\end{center}

\begin{verbatim}
a, b, c = 9, 8, 6
def h(a, b=7, c=7):
    print(a, b, c, end=" ")

h(b=5, a=2, c=4)
print(a, b, c)
\end{verbatim}



\begin{center}
\textbf{Завдання 9}
\end{center}

\begin{verbatim}
a,b,c=5,9,2
def g(b):
    a=2
    b+=4
    c=5
    return a+b+c

a,b,c=4,0,7
print(g(a),a,b,c)
\end{verbatim}



\begin{center}
\textbf{Завдання 10}
\end{center}

Написати функцiю f, що отримує 2 числа і повертає логічну ознаку того, що перше число менше за куб другого. Стиль запису умови також оцінюється.\par

\end{minipage}

\end {large}}
\newpage

{\Large{06.11.2025 Контрольна робота \No 1, варіант  \No \textbf { 7 }\par}}
{
\begin {large}
У завданнях 1-9 вказати, що буде виведено за виконання.

Повідомлення інтерпретатора про помилку позначати ERROR

\begin{minipage}[t]{7.5cm}

\begin{center}
\textbf{Завдання 1}
\end{center}

\begin{verbatim}
print(8/9*7*3)
\end{verbatim}



\begin{center}
\textbf{Завдання 2}
\end{center}

\begin{verbatim}
print(1**-4-2)
\end{verbatim}



\begin{center}
\textbf{Завдання 4}
\end{center}

\begin{verbatim}
print("7" < False)
\end{verbatim}



\begin{center}
\textbf{Завдання 6}
\end{center}

\begin{verbatim}

def f(n):
    print("f", end="");
    return n!=4

print(f(-5) or f(-9))
\end{verbatim}



\begin{center}
\textbf{Завдання 7}
\end{center}

\begin{verbatim}
def g(a,b):
    c=71
    if b<=-4:
        return 9
    elif a>4:
         return 6
    else: 
        c=5
    return c

print(g(-2,2))
\end{verbatim}

\end{minipage}
\vline \hspace{6mm}
\begin{minipage}[t]{10.5cm}

\begin{center}
\textbf{Завдання 3}
\end{center}

\begin{verbatim}
print(5.0<=False or not 4==8.0 or 8<7)
\end{verbatim}



\begin{center}
\textbf{Завдання 5}
\end{center}

\begin{verbatim}
print(3 is 5.0 >= 9 < 8.0)
\end{verbatim}



\begin{center}
\textbf{Завдання 8}
\end{center}

\begin{verbatim}
a, b, c = 8, 6, 9
def f(a, b, c):
    print(a, b, c, end=" ")

f(1, 0, a=5)
print(a, b, c)
\end{verbatim}



\begin{center}
\textbf{Завдання 9}
\end{center}

\begin{verbatim}
a,b,c=2,9,7
def h(a):
    a=5
    b-=2
    c=3
    return a+b+c

a,b,c=0,4,5
print(h(a),a,b,c)
\end{verbatim}



\begin{center}
\textbf{Завдання 10}
\end{center}

Написати функцiю f, що отримує число і повертає логічну ознаку того, що воно не належить вiдрiзку [2021; 2022]. Стиль запису умови також оцінюється.\par

\end{minipage}

\end {large}}
\newpage

{\Large{06.11.2025 Контрольна робота \No 1, варіант  \No \textbf { 8 }\par}}
{
\begin {large}
У завданнях 1-9 вказати, що буде виведено за виконання.

Повідомлення інтерпретатора про помилку позначати ERROR

\begin{minipage}[t]{7.5cm}

\begin{center}
\textbf{Завдання 1}
\end{center}

\begin{verbatim}
print(7//9%4*8)
\end{verbatim}



\begin{center}
\textbf{Завдання 2}
\end{center}

\begin{verbatim}
print(-1**1.0+True)
\end{verbatim}



\begin{center}
\textbf{Завдання 4}
\end{center}

\begin{verbatim}
print(6.0 <= "True")
\end{verbatim}



\begin{center}
\textbf{Завдання 6}
\end{center}

\begin{verbatim}

def f(n):
    print("f", end="");
    return n

print(f(8) and f(-3))
\end{verbatim}



\begin{center}
\textbf{Завдання 7}
\end{center}

\begin{verbatim}
def f(a,b):
    c=57
    if a:
        c=4
    elif a>=3:
         return 0
    else: 
        c=8
    return c

print(f(1,4))
\end{verbatim}

\end{minipage}
\vline \hspace{6mm}
\begin{minipage}[t]{10.5cm}

\begin{center}
\textbf{Завдання 3}
\end{center}

\begin{verbatim}
print(not 3>=5 and True>6 and 3.0!=9.0)
\end{verbatim}



\begin{center}
\textbf{Завдання 5}
\end{center}

\begin{verbatim}
print(8 is True != 5 > 6)
\end{verbatim}



\begin{center}
\textbf{Завдання 8}
\end{center}

\begin{verbatim}
a, b, c = 7, 6, 8
def h(a, b, c=9):
    print(a, b, c, end=" ")

h(1, b=3)
print(a, b, c)
\end{verbatim}



\begin{center}
\textbf{Завдання 9}
\end{center}

\begin{verbatim}
a,b,c=3,1,8
def g(a):
    a-=4
    b=1
    c=5
    return a+b+c

a,b,c=6,9,2
print(g(b),a,b,c)
\end{verbatim}



\begin{center}
\textbf{Завдання 10}
\end{center}

Написати функцiю f, що отримує рядок i повертає логічну ознаку того, що вiн складається рівно з однієї літери, яка є десятковою цифрою. Стиль запису умови також оцінюється.\par

\end{minipage}

\end {large}}
\newpage

{\Large{06.11.2025 Контрольна робота \No 1, варіант  \No \textbf { 9 }\par}}
{
\begin {large}
У завданнях 1-9 вказати, що буде виведено за виконання.

Повідомлення інтерпретатора про помилку позначати ERROR

\begin{minipage}[t]{7.5cm}

\begin{center}
\textbf{Завдання 1}
\end{center}

\begin{verbatim}
print(4/6%7*7)
\end{verbatim}



\begin{center}
\textbf{Завдання 2}
\end{center}

\begin{verbatim}
print(9**-1.0*-2)
\end{verbatim}



\begin{center}
\textbf{Завдання 4}
\end{center}

\begin{verbatim}
print(3 >= "5.0j")
\end{verbatim}



\begin{center}
\textbf{Завдання 6}
\end{center}

\begin{verbatim}

def f(n):
    print("f", end="");
    return n<0

print(f(4) and f(9))
\end{verbatim}



\begin{center}
\textbf{Завдання 7}
\end{center}

\begin{verbatim}
def f(c):
    u=72
    if c==-2: 
        u=6
    if c>5:
         u=8
    else:
         return 1
    return u

print(f(2))
\end{verbatim}

\end{minipage}
\vline \hspace{6mm}
\begin{minipage}[t]{10.5cm}

\begin{center}
\textbf{Завдання 3}
\end{center}

\begin{verbatim}
print(not 7.0<=6 and False>=9.0 and 5<9)
\end{verbatim}



\begin{center}
\textbf{Завдання 5}
\end{center}

\begin{verbatim}
print(False <= 9.0 == 4 < 3.0)
\end{verbatim}



\begin{center}
\textbf{Завдання 8}
\end{center}

\begin{verbatim}
a, b, c = 6, 7, 9
def g(a, b=8, c):
    print(a, b, c, end=" ")

g(2, 4, a=4)
print(a, b, c)
\end{verbatim}



\begin{center}
\textbf{Завдання 9}
\end{center}

\begin{verbatim}
a,b,c=8,3,4
def h(b):
    a*=1
    b=4
    c=3
    return a+b+c

a,b,c=7,1,5
print(h(a),a,b,c)
\end{verbatim}



\begin{center}
\textbf{Завдання 10}
\end{center}

Написати функцiю f, що отримує число і повертає логічну ознаку того, що воно не належить iнтервалу (2019; 2021). Стиль запису умови також оцінюється.\par

\end{minipage}

\end {large}}
\newpage

{\Large{06.11.2025 Контрольна робота \No 1, варіант  \No \textbf { 10 }\par}}
{
\begin {large}
У завданнях 1-9 вказати, що буде виведено за виконання.

Повідомлення інтерпретатора про помилку позначати ERROR

\begin{minipage}[t]{7.5cm}

\begin{center}
\textbf{Завдання 1}
\end{center}

\begin{verbatim}
print(8//8/9*3)
\end{verbatim}



\begin{center}
\textbf{Завдання 2}
\end{center}

\begin{verbatim}
print(4**1.0*False)
\end{verbatim}



\begin{center}
\textbf{Завдання 4}
\end{center}

\begin{verbatim}
print("2j" > 9)
\end{verbatim}



\begin{center}
\textbf{Завдання 6}
\end{center}

\begin{verbatim}

def f(n):
    print("f", end="");
    return n

print(f(-6) or f(-4))
\end{verbatim}



\begin{center}
\textbf{Завдання 7}
\end{center}

\begin{verbatim}
def f(a):
    x=38
    if a: 
        return 7
    elif a!=-5:
         x=9
    else:
         return 0
    return x

print(f(7))
\end{verbatim}

\end{minipage}
\vline \hspace{6mm}
\begin{minipage}[t]{10.5cm}

\begin{center}
\textbf{Завдання 3}
\end{center}

\begin{verbatim}
print(not 3.0>4.0 or True!=2 or 7==8)
\end{verbatim}



\begin{center}
\textbf{Завдання 5}
\end{center}

\begin{verbatim}
print(2.0 >= 8 is 4.0 > True)
\end{verbatim}



\begin{center}
\textbf{Завдання 8}
\end{center}

\begin{verbatim}
a, b, c = 6, 9, 7
def f(a, b=8, c=8):
    print(a, b, c, end=" ")

f(1, c=3, b=5)
print(a, b, c)
\end{verbatim}



\begin{center}
\textbf{Завдання 9}
\end{center}

\begin{verbatim}
a,b,c=7,8,0
def g(b):
    a=3
    b=2
    c=4
    return a+b+c

a,b,c=5,2,8
print(g(b),a,b,c)
\end{verbatim}



\begin{center}
\textbf{Завдання 10}
\end{center}

Написати функцiю f, що отримує число і повертає логічну ознаку того, що воно не належить вiдрiзку [2018; 2019]. Стиль запису умови також оцінюється.\par

\end{minipage}

\end {large}}
\newpage

{\Large{06.11.2025 Контрольна робота \No 1, варіант  \No \textbf { 11 }\par}}
{
\begin {large}
У завданнях 1-9 вказати, що буде виведено за виконання.

Повідомлення інтерпретатора про помилку позначати ERROR

\begin{minipage}[t]{7.5cm}

\begin{center}
\textbf{Завдання 1}
\end{center}

\begin{verbatim}
print(5//3*5*9)
\end{verbatim}



\begin{center}
\textbf{Завдання 2}
\end{center}

\begin{verbatim}
print(-3**4+1)
\end{verbatim}



\begin{center}
\textbf{Завдання 4}
\end{center}

\begin{verbatim}
print(False != "4.0")
\end{verbatim}



\begin{center}
\textbf{Завдання 6}
\end{center}

\begin{verbatim}

def f(n):
    print("f", end="");
    return n

print(f(-9) or f(-8))
\end{verbatim}



\begin{center}
\textbf{Завдання 7}
\end{center}

\begin{verbatim}
def g(a,b):
    c=41
    if a:
        return 3
    if b<1:
         return 2
    else: 
        c=5
    return c

print(g(0,1))
\end{verbatim}

\end{minipage}
\vline \hspace{6mm}
\begin{minipage}[t]{10.5cm}

\begin{center}
\textbf{Завдання 3}
\end{center}

\begin{verbatim}
print(not 4>3 and 6.0!=8 or True<8.0)
\end{verbatim}



\begin{center}
\textbf{Завдання 5}
\end{center}

\begin{verbatim}
print(5 < 7 <= 9 >= 2)
\end{verbatim}



\begin{center}
\textbf{Завдання 8}
\end{center}

\begin{verbatim}
a, b, c = 6, 7, 9
def g(a, b, c):
    print(a, b, c, end=" ")

g(a=2, 0, c=4)
print(a, b, c)
\end{verbatim}



\begin{center}
\textbf{Завдання 9}
\end{center}

\begin{verbatim}
a,b,c=4,7,6
def f(b):
    a=3
    b=5
    c=2
    return a+b+c

a,b,c=0,1,5
print(f(b),a,b,c)
\end{verbatim}



\begin{center}
\textbf{Завдання 10}
\end{center}

Написати функцiю f, що отримує число і повертає логічну ознаку того, що воно належить вiдрiзку [2021; 2022]. Стиль запису умови також оцінюється.\par

\end{minipage}

\end {large}}
\newpage

{\Large{06.11.2025 Контрольна робота \No 1, варіант  \No \textbf { 12 }\par}}
{
\begin {large}
У завданнях 1-9 вказати, що буде виведено за виконання.

Повідомлення інтерпретатора про помилку позначати ERROR

\begin{minipage}[t]{7.5cm}

\begin{center}
\textbf{Завдання 1}
\end{center}

\begin{verbatim}
print(6/4//3*5)
\end{verbatim}



\begin{center}
\textbf{Завдання 2}
\end{center}

\begin{verbatim}
print(4**-3--1)
\end{verbatim}



\begin{center}
\textbf{Завдання 4}
\end{center}

\begin{verbatim}
print(3.0 == "True")
\end{verbatim}



\begin{center}
\textbf{Завдання 6}
\end{center}

\begin{verbatim}

def f(n):
    print("f", end="");
    return n>-4

print(f(-5) and f(8))
\end{verbatim}



\begin{center}
\textbf{Завдання 7}
\end{center}

\begin{verbatim}
def g(d):
    z=18
    if d<-1: 
        z=1
    if d!=-4:
         z=3
    else:
         return 6
    return z

print(g(6))
\end{verbatim}

\end{minipage}
\vline \hspace{6mm}
\begin{minipage}[t]{10.5cm}

\begin{center}
\textbf{Завдання 3}
\end{center}

\begin{verbatim}
print(not 6<=2.0 or 4==2 and 5.0>=False)
\end{verbatim}



\begin{center}
\textbf{Завдання 5}
\end{center}

\begin{verbatim}
print(7.0 == False != 3 == 6.0)
\end{verbatim}



\begin{center}
\textbf{Завдання 8}
\end{center}

\begin{verbatim}
a, b, c = 9, 7, 6
def h(a, b, c=8):
    print(a, b, c, end=" ")

h(b=3, c=2, 1)
print(a, b, c)
\end{verbatim}



\begin{center}
\textbf{Завдання 9}
\end{center}

\begin{verbatim}
a,b,c=6,0,1
def f(a):
    a+=2
    b=2
    c=5
    return a+b+c

a,b,c=0,9,8
print(f(b),a,b,c)
\end{verbatim}



\begin{center}
\textbf{Завдання 10}
\end{center}

Написати функцiю f, що отримує 2 числа і повертає логічну ознаку того, що перше число безпечно дiлити на друге. Стиль запису умови також оцінюється.\par

\end{minipage}

\end {large}}
\newpage

{\Large{06.11.2025 Контрольна робота \No 1, варіант  \No \textbf { 13 }\par}}
{
\begin {large}
У завданнях 1-9 вказати, що буде виведено за виконання.

Повідомлення інтерпретатора про помилку позначати ERROR

\begin{minipage}[t]{7.5cm}

\begin{center}
\textbf{Завдання 1}
\end{center}

\begin{verbatim}
print(9//5//6*6)
\end{verbatim}



\begin{center}
\textbf{Завдання 2}
\end{center}

\begin{verbatim}
print(-1**-2-2)
\end{verbatim}



\begin{center}
\textbf{Завдання 4}
\end{center}

\begin{verbatim}
print("8.0j" >= False)
\end{verbatim}



\begin{center}
\textbf{Завдання 6}
\end{center}

\begin{verbatim}

def f(n):
    print("f", end="");
    return n

print(f(-1) and f(4))
\end{verbatim}



\begin{center}
\textbf{Завдання 7}
\end{center}

\begin{verbatim}
def h(a):
    u=80
    if a==-3: 
        u=5
    elif a<3:
         return 7
    else:
         u=8
    return u

print(h(4))
\end{verbatim}

\end{minipage}
\vline \hspace{6mm}
\begin{minipage}[t]{10.5cm}

\begin{center}
\textbf{Завдання 3}
\end{center}

\begin{verbatim}
print(not 7>=9.0 and 5<=True and 4.0>3)
\end{verbatim}



\begin{center}
\textbf{Завдання 5}
\end{center}

\begin{verbatim}
print(False != 7 < 2.0 is 6.0)
\end{verbatim}



\begin{center}
\textbf{Завдання 8}
\end{center}

\begin{verbatim}
a, b, c = 7, 8, 6
def f(a, b=9, c=7):
    print(a, b, c, end=" ")

f(5, b=0)
print(a, b, c)
\end{verbatim}



\begin{center}
\textbf{Завдання 9}
\end{center}

\begin{verbatim}
a,b,c=4,7,2
def f(b):
    a=4
    b*=1
    c=3
    return a+b+c

a,b,c=3,5,6
print(f(a),a,b,c)
\end{verbatim}



\begin{center}
\textbf{Завдання 10}
\end{center}

Написати функцiю f, що отримує 2 числа і повертає логічну ознаку того, що перше число бiльше за квадрат другого. Стиль запису умови також оцінюється.\par

\end{minipage}

\end {large}}
\newpage

{\Large{06.11.2025 Контрольна робота \No 1, варіант  \No \textbf { 14 }\par}}
{
\begin {large}
У завданнях 1-9 вказати, що буде виведено за виконання.

Повідомлення інтерпретатора про помилку позначати ERROR

\begin{minipage}[t]{7.5cm}

\begin{center}
\textbf{Завдання 1}
\end{center}

\begin{verbatim}
print(3/7*8*4)
\end{verbatim}



\begin{center}
\textbf{Завдання 2}
\end{center}

\begin{verbatim}
print(1**-0.5*False)
\end{verbatim}



\begin{center}
\textbf{Завдання 4}
\end{center}

\begin{verbatim}
print("True" < "8")
\end{verbatim}



\begin{center}
\textbf{Завдання 6}
\end{center}

\begin{verbatim}

def f(n):
    print("f", end="");
    return n!=1

print(f(2) or f(9))
\end{verbatim}



\begin{center}
\textbf{Завдання 7}
\end{center}

\begin{verbatim}
def h(a,b):
    c=51
    if b!=5:
        return 8
    if b==-2:
         c=0
    else: 
        return 1
    return c

print(h(6,5))
\end{verbatim}

\end{minipage}
\vline \hspace{6mm}
\begin{minipage}[t]{10.5cm}

\begin{center}
\textbf{Завдання 3}
\end{center}

\begin{verbatim}
print(not 2.0!=9 or 2<False or 9==3.0)
\end{verbatim}



\begin{center}
\textbf{Завдання 5}
\end{center}

\begin{verbatim}
print(9 > 2 <= 3.0 >= 3)
\end{verbatim}



\begin{center}
\textbf{Завдання 8}
\end{center}

\begin{verbatim}
a, b, c = 8, 6, 9
def h(a, b, c=8):
    print(a, b, c, end=" ")

h(4, 3, 5)
print(a, b, c)
\end{verbatim}



\begin{center}
\textbf{Завдання 9}
\end{center}

\begin{verbatim}
a,b,c=4,6,0
def h(a):
    a-=5
    b=3
    c=2
    return a+b+c

a,b,c=5,7,1
print(h(a),a,b,c)
\end{verbatim}



\begin{center}
\textbf{Завдання 10}
\end{center}

Написати функцiю f, що отримує число і повертає логічну ознаку того, що воно належить iнтервалу (2020; 2021). Стиль запису умови також оцінюється.\par

\end{minipage}

\end {large}}
\newpage

{\Large{06.11.2025 Контрольна робота \No 1, варіант  \No \textbf { 15 }\par}}
{
\begin {large}
У завданнях 1-9 вказати, що буде виведено за виконання.

Повідомлення інтерпретатора про помилку позначати ERROR

\begin{minipage}[t]{7.5cm}

\begin{center}
\textbf{Завдання 1}
\end{center}

\begin{verbatim}
print(6//6%7*7)
\end{verbatim}



\begin{center}
\textbf{Завдання 2}
\end{center}

\begin{verbatim}
print(3**False+1)
\end{verbatim}



\begin{center}
\textbf{Завдання 4}
\end{center}

\begin{verbatim}
print(8 != 5.0)
\end{verbatim}



\begin{center}
\textbf{Завдання 6}
\end{center}

\begin{verbatim}

def f(n):
    print("f", end="");
    return n>-2

print(f(-6) and f(3))
\end{verbatim}



\begin{center}
\textbf{Завдання 7}
\end{center}

\begin{verbatim}
def f(a,b):
    c=54
    if a>2:
        c=3
    elif a<=-5:
         return 9
    else: 
        c=7
    return c

print(f(-7,0))
\end{verbatim}

\end{minipage}
\vline \hspace{6mm}
\begin{minipage}[t]{10.5cm}

\begin{center}
\textbf{Завдання 3}
\end{center}

\begin{verbatim}
print(5!=8.0 and not 7.0>7 and 6>=False)
\end{verbatim}



\begin{center}
\textbf{Завдання 5}
\end{center}

\begin{verbatim}
print(5 > 5.0 is 6 != True)
\end{verbatim}



\begin{center}
\textbf{Завдання 8}
\end{center}

\begin{verbatim}
a, b, c = 7, 6, 9
def g(a, b, c):
    print(a, b, c, end=" ")

g(0, 1)
print(a, b, c)
\end{verbatim}



\begin{center}
\textbf{Завдання 9}
\end{center}

\begin{verbatim}
a,b,c=3,2,8
def g(b):
    a+=1
    b=4
    c=2
    return a+b+c

a,b,c=9,5,8
print(g(b),a,b,c)
\end{verbatim}



\begin{center}
\textbf{Завдання 10}
\end{center}

Написати функцiю f, що отримує 2 цілих числа i повертає логічну ознаку того, що їх дванадцятковий запис закінчується на одну ту саму цифру. Стиль запису умови також оцінюється.\par

\end{minipage}

\end {large}}
\newpage

{\Large{06.11.2025 Контрольна робота \No 1, варіант  \No \textbf { 16 }\par}}
{
\begin {large}
У завданнях 1-9 вказати, що буде виведено за виконання.

Повідомлення інтерпретатора про помилку позначати ERROR

\begin{minipage}[t]{7.5cm}

\begin{center}
\textbf{Завдання 1}
\end{center}

\begin{verbatim}
print(3/7/9*9)
\end{verbatim}



\begin{center}
\textbf{Завдання 2}
\end{center}

\begin{verbatim}
print(9**0.5*-1)
\end{verbatim}



\begin{center}
\textbf{Завдання 4}
\end{center}

\begin{verbatim}
print(False <= "9.0")
\end{verbatim}



\begin{center}
\textbf{Завдання 6}
\end{center}

\begin{verbatim}

def f(n):
    print("f", end="");
    return n

print(f(-3) or f(7))
\end{verbatim}



\begin{center}
\textbf{Завдання 7}
\end{center}

\begin{verbatim}
def f(d):
    v=76
    if d<=5: 
        return 9
    if d>1:
         return 4
    else:
         v=2
    return v

print(f(9))
\end{verbatim}

\end{minipage}
\vline \hspace{6mm}
\begin{minipage}[t]{10.5cm}

\begin{center}
\textbf{Завдання 3}
\end{center}

\begin{verbatim}
print(not 5.0<True or 8<=4 or 3==6.0)
\end{verbatim}



\begin{center}
\textbf{Завдання 5}
\end{center}

\begin{verbatim}
print(False <= 9.0 < 4 == 7.0)
\end{verbatim}



\begin{center}
\textbf{Завдання 8}
\end{center}

\begin{verbatim}
a, b, c = 7, 9, 8
def f(a, b=6, c):
    print(a, b, c, end=" ")

f(a=2, c=5, b=2)
print(a, b, c)
\end{verbatim}



\begin{center}
\textbf{Завдання 9}
\end{center}

\begin{verbatim}
a,b,c=3,9,7
def f(a):
    a=1
    b=4
    c=3
    return a+b+c

a,b,c=4,1,5
print(f(a),a,b,c)
\end{verbatim}



\begin{center}
\textbf{Завдання 10}
\end{center}

Написати функцiю f, що отримує 3 числа i повертає логічну ознаку того, що вони попарно рiвнi. Стиль запису умови також оцінюється.\par

\end{minipage}

\end {large}}
\newpage

{\Large{06.11.2025 Контрольна робота \No 1, варіант  \No \textbf { 17 }\par}}
{
\begin {large}
У завданнях 1-9 вказати, що буде виведено за виконання.

Повідомлення інтерпретатора про помилку позначати ERROR

\begin{minipage}[t]{7.5cm}

\begin{center}
\textbf{Завдання 1}
\end{center}

\begin{verbatim}
print(4//5/3*5)
\end{verbatim}



\begin{center}
\textbf{Завдання 2}
\end{center}

\begin{verbatim}
print(4**0.5*True)
\end{verbatim}



\begin{center}
\textbf{Завдання 4}
\end{center}

\begin{verbatim}
print("True" == "1j")
\end{verbatim}



\begin{center}
\textbf{Завдання 6}
\end{center}

\begin{verbatim}

def f(n):
    print("f", end="");
    return n<=3

print(f(0) or f(5))
\end{verbatim}



\begin{center}
\textbf{Завдання 7}
\end{center}

\begin{verbatim}
def f(a,b):
    c=93
    if b<-3:
        c=6
    if a>=-1:
         return 4
    else: 
        return 9
    return c

print(f(2,6))
\end{verbatim}

\end{minipage}
\vline \hspace{6mm}
\begin{minipage}[t]{10.5cm}

\begin{center}
\textbf{Завдання 3}
\end{center}

\begin{verbatim}
print(not 9!=9.0 and True<5 or 7.0==7)
\end{verbatim}



\begin{center}
\textbf{Завдання 5}
\end{center}

\begin{verbatim}
print(8 >= 4 is True < 2)
\end{verbatim}



\begin{center}
\textbf{Завдання 8}
\end{center}

\begin{verbatim}
a, b, c = 7, 8, 6
def h(a, b, c=9):
    print(a, b, c, end=" ")

h(3, a=4)
print(a, b, c)
\end{verbatim}



\begin{center}
\textbf{Завдання 9}
\end{center}

\begin{verbatim}
a,b,c=1,6,0
def g(a):
    a*=4
    b=1
    c=3
    return a+b+c

a,b,c=3,9,2
print(g(b),a,b,c)
\end{verbatim}



\begin{center}
\textbf{Завдання 10}
\end{center}

Написати функцiю f, що отримує рядок i повертає логічну ознаку того, що вiн складається рівно з однієї малої латинської літери. Стиль запису умови також оцінюється.\par

\end{minipage}

\end {large}}
\newpage

{\Large{06.11.2025 Контрольна робота \No 1, варіант  \No \textbf { 18 }\par}}
{
\begin {large}
У завданнях 1-9 вказати, що буде виведено за виконання.

Повідомлення інтерпретатора про помилку позначати ERROR

\begin{minipage}[t]{7.5cm}

\begin{center}
\textbf{Завдання 1}
\end{center}

\begin{verbatim}
print(7/4*8*8)
\end{verbatim}



\begin{center}
\textbf{Завдання 2}
\end{center}

\begin{verbatim}
print(-2**-2+-2)
\end{verbatim}



\begin{center}
\textbf{Завдання 4}
\end{center}

\begin{verbatim}
print(3 > 6.0)
\end{verbatim}



\begin{center}
\textbf{Завдання 6}
\end{center}

\begin{verbatim}

def f(n):
    print("f", end="");
    return n

print(f(6) and f(-2))
\end{verbatim}



\begin{center}
\textbf{Завдання 7}
\end{center}

\begin{verbatim}
def h(a,b):
    c=12
    if a!=2:
        c=2
    elif b>5:
         c=1
    else: 
        return 4
    return c

print(h(9,-3))
\end{verbatim}

\end{minipage}
\vline \hspace{6mm}
\begin{minipage}[t]{10.5cm}

\begin{center}
\textbf{Завдання 3}
\end{center}

\begin{verbatim}
print(2<=False or not 5.0>=8 and 4>4.0)
\end{verbatim}



\begin{center}
\textbf{Завдання 5}
\end{center}

\begin{verbatim}
print(8.0 == 4.0 > 8 >= 3)
\end{verbatim}



\begin{center}
\textbf{Завдання 8}
\end{center}

\begin{verbatim}
a, b, c = 6, 9, 8
def g(a, b, c):
    print(a, b, c, end=" ")

g(a=0, 1, b=4)
print(a, b, c)
\end{verbatim}



\begin{center}
\textbf{Завдання 9}
\end{center}

\begin{verbatim}
a,b,c=4,7,3
def f(b):
    a+=5
    b=2
    c=4
    return a+b+c

a,b,c=8,4,7
print(f(a),a,b,c)
\end{verbatim}



\begin{center}
\textbf{Завдання 10}
\end{center}

Написати функцiю f, що отримує 2 цілих числа i повертає логічну ознаку того, що їх сімковий запис закінчується на одну ту саму цифру. Стиль запису умови також оцінюється.\par

\end{minipage}

\end {large}}
\newpage

{\Large{06.11.2025 Контрольна робота \No 1, варіант  \No \textbf { 19 }\par}}
{
\begin {large}
У завданнях 1-9 вказати, що буде виведено за виконання.

Повідомлення інтерпретатора про помилку позначати ERROR

\begin{minipage}[t]{7.5cm}

\begin{center}
\textbf{Завдання 1}
\end{center}

\begin{verbatim}
print(9/9//6*3)
\end{verbatim}



\begin{center}
\textbf{Завдання 2}
\end{center}

\begin{verbatim}
print(4**1.0*False)
\end{verbatim}



\begin{center}
\textbf{Завдання 4}
\end{center}

\begin{verbatim}
print("7.0" > 7)
\end{verbatim}



\begin{center}
\textbf{Завдання 6}
\end{center}

\begin{verbatim}

def f(n):
    print("f", end="");
    return n==-4

print(f(-7) or f(1))
\end{verbatim}



\begin{center}
\textbf{Завдання 7}
\end{center}

\begin{verbatim}
def g(a,b):
    c=47
    if a:
        c=8
    elif b<=0:
         return 7
    else: 
        return 6
    return c

print(g(5,-6))
\end{verbatim}

\end{minipage}
\vline \hspace{6mm}
\begin{minipage}[t]{10.5cm}

\begin{center}
\textbf{Завдання 3}
\end{center}

\begin{verbatim}
print(not 6!=False or 3>9 and 3.0<2.0)
\end{verbatim}



\begin{center}
\textbf{Завдання 5}
\end{center}

\begin{verbatim}
print(9 <= 2.0 != 3.0 < 6)
\end{verbatim}



\begin{center}
\textbf{Завдання 8}
\end{center}

\begin{verbatim}
a, b, c = 7, 9, 6
def h(a, b=8, c=7):
    print(a, b, c, end=" ")

h(1, 5, c=3)
print(a, b, c)
\end{verbatim}



\begin{center}
\textbf{Завдання 9}
\end{center}

\begin{verbatim}
a,b,c=2,1,6
def h(a):
    a-=3
    b=1
    c=5
    return a+b+c

a,b,c=5,9,0
print(h(b),a,b,c)
\end{verbatim}



\begin{center}
\textbf{Завдання 10}
\end{center}

Написати функцiю f, що отримує 2 числа і повертає логічну ознаку того, що перше число менше за квадрат другого. Стиль запису умови також оцінюється.\par

\end{minipage}

\end {large}}
\newpage

{\Large{06.11.2025 Контрольна робота \No 1, варіант  \No \textbf { 20 }\par}}
{
\begin {large}
У завданнях 1-9 вказати, що буде виведено за виконання.

Повідомлення інтерпретатора про помилку позначати ERROR

\begin{minipage}[t]{7.5cm}

\begin{center}
\textbf{Завдання 1}
\end{center}

\begin{verbatim}
print(8//8%5*4)
\end{verbatim}



\begin{center}
\textbf{Завдання 2}
\end{center}

\begin{verbatim}
print(3**-1--1)
\end{verbatim}



\begin{center}
\textbf{Завдання 4}
\end{center}

\begin{verbatim}
print("False" <= "2")
\end{verbatim}



\begin{center}
\textbf{Завдання 6}
\end{center}

\begin{verbatim}

def f(n):
    print("f", end="");
    return n

print(f(7) and f(9))
\end{verbatim}



\begin{center}
\textbf{Завдання 7}
\end{center}

\begin{verbatim}
def f(a,b):
    c=62
    if a<-5:
        c=3
    if b==1:
         return 5
    else: 
        c=0
    return c

print(f(8,-5))
\end{verbatim}

\end{minipage}
\vline \hspace{6mm}
\begin{minipage}[t]{10.5cm}

\begin{center}
\textbf{Завдання 3}
\end{center}

\begin{verbatim}
print(not 5==True and 6<=8.0 or 8>=6.0)
\end{verbatim}



\begin{center}
\textbf{Завдання 5}
\end{center}

\begin{verbatim}
print(False is 5 > True <= 7.0)
\end{verbatim}



\begin{center}
\textbf{Завдання 8}
\end{center}

\begin{verbatim}
a, b, c = 7, 6, 9
def g(a, b=8, c):
    print(a, b, c, end=" ")

g(2, 0)
print(a, b, c)
\end{verbatim}



\begin{center}
\textbf{Завдання 9}
\end{center}

\begin{verbatim}
a,b,c=3,6,8
def h(b):
    a=1
    b=2
    c=5
    return a+b+c

a,b,c=2,0,9
print(h(b),a,b,c)
\end{verbatim}



\begin{center}
\textbf{Завдання 10}
\end{center}

Написати функцiю f, що отримує число і повертає логічну ознаку того, що воно належить вiдрiзку [2016; 2017]. Стиль запису умови також оцінюється.\par

\end{minipage}

\end {large}}
\newpage

{\Large{06.11.2025 Контрольна робота \No 1, варіант  \No \textbf { 21 }\par}}
{
\begin {large}
У завданнях 1-9 вказати, що буде виведено за виконання.

Повідомлення інтерпретатора про помилку позначати ERROR

\begin{minipage}[t]{7.5cm}

\begin{center}
\textbf{Завдання 1}
\end{center}

\begin{verbatim}
print(5//3/4*6)
\end{verbatim}



\begin{center}
\textbf{Завдання 2}
\end{center}

\begin{verbatim}
print(4**4-1)
\end{verbatim}



\begin{center}
\textbf{Завдання 4}
\end{center}

\begin{verbatim}
print(3.0 >= True)
\end{verbatim}



\begin{center}
\textbf{Завдання 6}
\end{center}

\begin{verbatim}

def f(n):
    print("f", end="");
    return n<-3

print(f(-2) or f(5))
\end{verbatim}



\begin{center}
\textbf{Завдання 7}
\end{center}

\begin{verbatim}
def h(c):
    w=69
    if c: 
        return 6
    elif c>=4:
         w=2
    else:
         w=4
    return w

print(h(-9))
\end{verbatim}

\end{minipage}
\vline \hspace{6mm}
\begin{minipage}[t]{10.5cm}

\begin{center}
\textbf{Завдання 3}
\end{center}

\begin{verbatim}
print(not 6.0==False or 3>7 and 2.0<=4)
\end{verbatim}



\begin{center}
\textbf{Завдання 5}
\end{center}

\begin{verbatim}
print(7 == 7 >= 6.0 != 4)
\end{verbatim}



\begin{center}
\textbf{Завдання 8}
\end{center}

\begin{verbatim}
a, b, c = 8, 9, 7
def f(a, b=6, c=8):
    print(a, b, c, end=" ")

f(5, 3, 4)
print(a, b, c)
\end{verbatim}



\begin{center}
\textbf{Завдання 9}
\end{center}

\begin{verbatim}
a,b,c=7,9,2
def h(b):
    a=3
    b-=4
    c=5
    return a+b+c

a,b,c=4,0,5
print(h(a),a,b,c)
\end{verbatim}



\begin{center}
\textbf{Завдання 10}
\end{center}

Написати функцiю f, що отримує 2 числа і повертає логічну ознаку того, що перше число бiльше за куб другого. Стиль запису умови також оцінюється.\par

\end{minipage}

\end {large}}
\newpage

{\Large{06.11.2025 Контрольна робота \No 1, варіант  \No \textbf { 22 }\par}}
{
\begin {large}
У завданнях 1-9 вказати, що буде виведено за виконання.

Повідомлення інтерпретатора про помилку позначати ERROR

\begin{minipage}[t]{7.5cm}

\begin{center}
\textbf{Завдання 1}
\end{center}

\begin{verbatim}
print(4/5*9*8)
\end{verbatim}



\begin{center}
\textbf{Завдання 2}
\end{center}

\begin{verbatim}
print(9**-0.5*True)
\end{verbatim}



\begin{center}
\textbf{Завдання 4}
\end{center}

\begin{verbatim}
print(False < 4)
\end{verbatim}



\begin{center}
\textbf{Завдання 6}
\end{center}

\begin{verbatim}

def f(n):
    print("f", end="");
    return n

print(f(-4) and f(4))
\end{verbatim}



\begin{center}
\textbf{Завдання 7}
\end{center}

\begin{verbatim}
def h(a,b):
    c=55
    if a:
        c=8
    if a>=-2:
         return 0
    else: 
        c=2
    return c

print(h(-9,-7))
\end{verbatim}

\end{minipage}
\vline \hspace{6mm}
\begin{minipage}[t]{10.5cm}

\begin{center}
\textbf{Завдання 3}
\end{center}

\begin{verbatim}
print(2>=True and not 8.0!=4 or 3<3.0)
\end{verbatim}



\begin{center}
\textbf{Завдання 5}
\end{center}

\begin{verbatim}
print(8 is True < 4.0 <= 9)
\end{verbatim}



\begin{center}
\textbf{Завдання 8}
\end{center}

\begin{verbatim}
a, b, c = 7, 9, 6
def f(a, b=8, c):
    print(a, b, c, end=" ")

f(2, c=0, b=1)
print(a, b, c)
\end{verbatim}



\begin{center}
\textbf{Завдання 9}
\end{center}

\begin{verbatim}
a,b,c=8,6,3
def f(b):
    a=1
    b*=2
    c=5
    return a+b+c

a,b,c=1,9,6
print(f(b),a,b,c)
\end{verbatim}



\begin{center}
\textbf{Завдання 10}
\end{center}

Написати функцiю f, що отримує рядок i повертає логічну ознаку того, що вiн складається рівно з однієї літери, яка є вісімковою цифрою. Стиль запису умови також оцінюється.\par

\end{minipage}

\end {large}}
\newpage

{\Large{06.11.2025 Контрольна робота \No 1, варіант  \No \textbf { 23 }\par}}
{
\begin {large}
У завданнях 1-9 вказати, що буде виведено за виконання.

Повідомлення інтерпретатора про помилку позначати ERROR

\begin{minipage}[t]{7.5cm}

\begin{center}
\textbf{Завдання 1}
\end{center}

\begin{verbatim}
print(8/4//8*9)
\end{verbatim}



\begin{center}
\textbf{Завдання 2}
\end{center}

\begin{verbatim}
print(-2**2+2)
\end{verbatim}



\begin{center}
\textbf{Завдання 4}
\end{center}

\begin{verbatim}
print(4.0 != "True")
\end{verbatim}



\begin{center}
\textbf{Завдання 6}
\end{center}

\begin{verbatim}

def f(n):
    print("f", end="");
    return n>=2

print(f(3) and f(-6))
\end{verbatim}



\begin{center}
\textbf{Завдання 7}
\end{center}

\begin{verbatim}
def g(a,b):
    c=28
    if b<-3:
        return 4
    elif b<=-4:
         return 3
    else: 
        c=1
    return c

print(g(-4,7))
\end{verbatim}

\end{minipage}
\vline \hspace{6mm}
\begin{minipage}[t]{10.5cm}

\begin{center}
\textbf{Завдання 3}
\end{center}

\begin{verbatim}
print(5.0<2 or not 9==8 or 4.0>=False)
\end{verbatim}



\begin{center}
\textbf{Завдання 5}
\end{center}

\begin{verbatim}
print(6 == 9.0 != 3 > False)
\end{verbatim}



\begin{center}
\textbf{Завдання 8}
\end{center}

\begin{verbatim}
a, b, c = 6, 9, 7
def h(a, b, c):
    print(a, b, c, end=" ")

h(a=2, b=3, c=5)
print(a, b, c)
\end{verbatim}



\begin{center}
\textbf{Завдання 9}
\end{center}

\begin{verbatim}
a,b,c=8,1,0
def g(a):
    a=4
    b+=3
    c=1
    return a+b+c

a,b,c=7,2,5
print(g(a),a,b,c)
\end{verbatim}



\begin{center}
\textbf{Завдання 10}
\end{center}

Написати функцiю f, що отримує 2 числа і повертає логічну ознаку того, що друге більше. Стиль запису умови також оцінюється.\par

\end{minipage}

\end {large}}
\newpage

{\Large{06.11.2025 Контрольна робота \No 1, варіант  \No \textbf { 24 }\par}}
{
\begin {large}
У завданнях 1-9 вказати, що буде виведено за виконання.

Повідомлення інтерпретатора про помилку позначати ERROR

\begin{minipage}[t]{7.5cm}

\begin{center}
\textbf{Завдання 1}
\end{center}

\begin{verbatim}
print(3//7%7*4)
\end{verbatim}



\begin{center}
\textbf{Завдання 2}
\end{center}

\begin{verbatim}
print(1**-1.0*-2)
\end{verbatim}



\begin{center}
\textbf{Завдання 4}
\end{center}

\begin{verbatim}
print("8.0j" == "9j")
\end{verbatim}



\begin{center}
\textbf{Завдання 6}
\end{center}

\begin{verbatim}

def f(n):
    print("f", end="");
    return n

print(f(-5) or f(-9))
\end{verbatim}



\begin{center}
\textbf{Завдання 7}
\end{center}

\begin{verbatim}
def g(a,b):
    c=33
    if b>4:
        return 9
    if b>=3:
         c=7
    else: 
        return 5
    return c

print(g(-5,-2))
\end{verbatim}

\end{minipage}
\vline \hspace{6mm}
\begin{minipage}[t]{10.5cm}

\begin{center}
\textbf{Завдання 3}
\end{center}

\begin{verbatim}
print(not 7!=9.0 and 7.0>True and 5<=6)
\end{verbatim}



\begin{center}
\textbf{Завдання 5}
\end{center}

\begin{verbatim}
print(5 >= 8.0 >= 5.0 is 2)
\end{verbatim}



\begin{center}
\textbf{Завдання 8}
\end{center}

\begin{verbatim}
a, b, c = 9, 7, 6
def g(a, b, c=8):
    print(a, b, c, end=" ")

g(b=1, c=0, 4)
print(a, b, c)
\end{verbatim}



\begin{center}
\textbf{Завдання 9}
\end{center}

\begin{verbatim}
a,b,c=1,5,9
def g(a):
    a=4
    b=5
    c=2
    return a+b+c

a,b,c=0,7,3
print(g(a),a,b,c)
\end{verbatim}



\begin{center}
\textbf{Завдання 10}
\end{center}

Написати функцiю f, що отримує число і повертає логічну ознаку того, що воно не належить iнтервалу (2017; 2018). Стиль запису умови також оцінюється.\par

\end{minipage}

\end {large}}
\newpage

{\Large{06.11.2025 Контрольна робота \No 1, варіант  \No \textbf { 25 }\par}}
{
\begin {large}
У завданнях 1-9 вказати, що буде виведено за виконання.

Повідомлення інтерпретатора про помилку позначати ERROR

\begin{minipage}[t]{7.5cm}

\begin{center}
\textbf{Завдання 1}
\end{center}

\begin{verbatim}
print(7/9//3*3)
\end{verbatim}



\begin{center}
\textbf{Завдання 2}
\end{center}

\begin{verbatim}
print(4**-1.0*2)
\end{verbatim}



\begin{center}
\textbf{Завдання 4}
\end{center}

\begin{verbatim}
print(5 != "6j")
\end{verbatim}



\begin{center}
\textbf{Завдання 6}
\end{center}

\begin{verbatim}

def f(n):
    print("f", end="");
    return n

print(f(-8) and f(-1))
\end{verbatim}



\begin{center}
\textbf{Завдання 7}
\end{center}

\begin{verbatim}
def g(b):
    y=92
    if b!=2: 
        return 7
    if b<=-2:
         y=8
    else:
         return 1
    return y

print(g(-1))
\end{verbatim}

\end{minipage}
\vline \hspace{6mm}
\begin{minipage}[t]{10.5cm}

\begin{center}
\textbf{Завдання 3}
\end{center}

\begin{verbatim}
print(not 7==True or 2.0>=8.0 and 9<=6)
\end{verbatim}



\begin{center}
\textbf{Завдання 5}
\end{center}

\begin{verbatim}
print(3.0 < 9 > 4 <= 2)
\end{verbatim}



\begin{center}
\textbf{Завдання 8}
\end{center}

\begin{verbatim}
a, b, c = 6, 8, 9
def g(a, b, c):
    print(a, b, c, end=" ")

g(0, 4, b=2)
print(a, b, c)
\end{verbatim}



\begin{center}
\textbf{Завдання 9}
\end{center}

\begin{verbatim}
a,b,c=3,4,3
def h(b):
    a=2
    b+=4
    c=3
    return a+b+c

a,b,c=5,9,8
print(h(a),a,b,c)
\end{verbatim}



\begin{center}
\textbf{Завдання 10}
\end{center}

Написати функцiю f, що отримує число і повертає логічну ознаку того, що воно не належить вiдрiзку [2019; 2020]. Стиль запису умови також оцінюється.\par

\end{minipage}

\end {large}}
\newpage

{\Large{06.11.2025 Контрольна робота \No 1, варіант  \No \textbf { 26 }\par}}
{
\begin {large}
У завданнях 1-9 вказати, що буде виведено за виконання.

Повідомлення інтерпретатора про помилку позначати ERROR

\begin{minipage}[t]{7.5cm}

\begin{center}
\textbf{Завдання 1}
\end{center}

\begin{verbatim}
print(9//8*4*7)
\end{verbatim}



\begin{center}
\textbf{Завдання 2}
\end{center}

\begin{verbatim}
print(9**-0.5*-1)
\end{verbatim}



\begin{center}
\textbf{Завдання 4}
\end{center}

\begin{verbatim}
print(False == "True")
\end{verbatim}



\begin{center}
\textbf{Завдання 6}
\end{center}

\begin{verbatim}

def f(n):
    print("f", end="");
    return n!=-1

print(f(8) or f(0))
\end{verbatim}



\begin{center}
\textbf{Завдання 7}
\end{center}

\begin{verbatim}
def h(c):
    z=86
    if c: 
        z=3
    elif c==-5:
         return 5
    else:
         z=0
    return z

print(h(-2))
\end{verbatim}

\end{minipage}
\vline \hspace{6mm}
\begin{minipage}[t]{10.5cm}

\begin{center}
\textbf{Завдання 3}
\end{center}

\begin{verbatim}
print(4.0>2 and not 7.0!=False or 3<4)
\end{verbatim}



\begin{center}
\textbf{Завдання 5}
\end{center}

\begin{verbatim}
print(5.0 == False != True == 7)
\end{verbatim}



\begin{center}
\textbf{Завдання 8}
\end{center}

\begin{verbatim}
a, b, c = 7, 6, 9
def h(a, b=8, c):
    print(a, b, c, end=" ")

h(a=3, 1, c=5)
print(a, b, c)
\end{verbatim}



\begin{center}
\textbf{Завдання 9}
\end{center}

\begin{verbatim}
a,b,c=2,6,4
def g(a):
    a*=1
    b=2
    c=5
    return a+b+c

a,b,c=1,7,0
print(g(b),a,b,c)
\end{verbatim}



\begin{center}
\textbf{Завдання 10}
\end{center}

Написати функцiю f, що отримує число і повертає логічну ознаку того, що воно належить iнтервалу (2016; 2017). Стиль запису умови також оцінюється.\par

\end{minipage}

\end {large}}
\newpage

{\Large{06.11.2025 Контрольна робота \No 1, варіант  \No \textbf { 27 }\par}}
{
\begin {large}
У завданнях 1-9 вказати, що буде виведено за виконання.

Повідомлення інтерпретатора про помилку позначати ERROR

\begin{minipage}[t]{7.5cm}

\begin{center}
\textbf{Завдання 1}
\end{center}

\begin{verbatim}
print(5/3/6*5)
\end{verbatim}



\begin{center}
\textbf{Завдання 2}
\end{center}

\begin{verbatim}
print(3**-1--2)
\end{verbatim}



\begin{center}
\textbf{Завдання 4}
\end{center}

\begin{verbatim}
print(1.0 <= "2.0")
\end{verbatim}



\begin{center}
\textbf{Завдання 6}
\end{center}

\begin{verbatim}

def f(n):
    print("f", end="");
    return n==3

print(f(1) and f(-7))
\end{verbatim}



\begin{center}
\textbf{Завдання 7}
\end{center}

\begin{verbatim}
def f(a,b):
    c=43
    if b!=-1:
        c=6
    elif a==2:
         c=1
    else: 
        return 4
    return c

print(f(7,3))
\end{verbatim}

\end{minipage}
\vline \hspace{6mm}
\begin{minipage}[t]{10.5cm}

\begin{center}
\textbf{Завдання 3}
\end{center}

\begin{verbatim}
print(not 5.0==5 and 8>2 and True!=9.0)
\end{verbatim}



\begin{center}
\textbf{Завдання 5}
\end{center}

\begin{verbatim}
print(9.0 < 7.0 is 5 >= 6)
\end{verbatim}



\begin{center}
\textbf{Завдання 8}
\end{center}

\begin{verbatim}
a, b, c = 7, 9, 7
def f(a, b=8, c=6):
    print(a, b, c, end=" ")

f(3, 0)
print(a, b, c)
\end{verbatim}



\begin{center}
\textbf{Завдання 9}
\end{center}

\begin{verbatim}
a,b,c=2,4,8
def f(b):
    a=1
    b=3
    c=1
    return a+b+c

a,b,c=6,7,2
print(f(a),a,b,c)
\end{verbatim}



\begin{center}
\textbf{Завдання 10}
\end{center}

Написати функцiю f, що отримує 2 цілих числа i повертає логічну ознаку того, що їх дев'ятковий запис закінчується на одну ту саму цифру. Стиль запису умови також оцінюється.\par

\end{minipage}

\end {large}}
\newpage

{\Large{06.11.2025 Контрольна робота \No 1, варіант  \No \textbf { 28 }\par}}
{
\begin {large}
У завданнях 1-9 вказати, що буде виведено за виконання.

Повідомлення інтерпретатора про помилку позначати ERROR

\begin{minipage}[t]{7.5cm}

\begin{center}
\textbf{Завдання 1}
\end{center}

\begin{verbatim}
print(6//6%5*6)
\end{verbatim}



\begin{center}
\textbf{Завдання 2}
\end{center}

\begin{verbatim}
print(4**2.0+1)
\end{verbatim}



\begin{center}
\textbf{Завдання 4}
\end{center}

\begin{verbatim}
print(True >= 5.0)
\end{verbatim}



\begin{center}
\textbf{Завдання 6}
\end{center}

\begin{verbatim}

def f(n):
    print("f", end="");
    return n

print(f(6) or f(2))
\end{verbatim}



\begin{center}
\textbf{Завдання 7}
\end{center}

\begin{verbatim}
def f(b):
    v=48
    if b<0: 
        return 9
    if b>5:
         return 1
    else:
         v=6
    return v

print(f(8))
\end{verbatim}

\end{minipage}
\vline \hspace{6mm}
\begin{minipage}[t]{10.5cm}

\begin{center}
\textbf{Завдання 3}
\end{center}

\begin{verbatim}
print(not 8>=7 or 6.0<=6 or 3.0<False)
\end{verbatim}



\begin{center}
\textbf{Завдання 5}
\end{center}

\begin{verbatim}
print(3 != 8 > 6.0 <= 5)
\end{verbatim}



\begin{center}
\textbf{Завдання 8}
\end{center}

\begin{verbatim}
a, b, c = 9, 7, 8
def h(a, b, c=6):
    print(a, b, c, end=" ")

h(2, a=4)
print(a, b, c)
\end{verbatim}



\begin{center}
\textbf{Завдання 9}
\end{center}

\begin{verbatim}
a,b,c=9,6,1
def h(a):
    a=4
    b=2
    c=5
    return a+b+c

a,b,c=4,5,8
print(h(b),a,b,c)
\end{verbatim}



\begin{center}
\textbf{Завдання 10}
\end{center}

Написати функцiю f, що отримує 2 числа і повертає логічну ознаку того, що хоча б одне з них нуль. Стиль запису умови також оцінюється.\par

\end{minipage}

\end {large}}
\newpage

{\Large{06.11.2025 Контрольна робота \No 1, варіант  \No \textbf { 29 }\par}}
{
\begin {large}
У завданнях 1-9 вказати, що буде виведено за виконання.

Повідомлення інтерпретатора про помилку позначати ERROR

\begin{minipage}[t]{7.5cm}

\begin{center}
\textbf{Завдання 1}
\end{center}

\begin{verbatim}
print(7//6*8*8)
\end{verbatim}



\begin{center}
\textbf{Завдання 2}
\end{center}

\begin{verbatim}
print(-3**-4-True)
\end{verbatim}



\begin{center}
\textbf{Завдання 4}
\end{center}

\begin{verbatim}
print("9.0j" < "3")
\end{verbatim}



\begin{center}
\textbf{Завдання 6}
\end{center}

\begin{verbatim}

def f(n):
    print("f", end="");
    return n

print(f(-3) or f(-6))
\end{verbatim}



\begin{center}
\textbf{Завдання 7}
\end{center}

\begin{verbatim}
def g(d):
    w=65
    if d>=-5: 
        return 0
    if d<1:
         w=2
    else:
         return 8
    return w

print(g(-6))
\end{verbatim}

\end{minipage}
\vline \hspace{6mm}
\begin{minipage}[t]{10.5cm}

\begin{center}
\textbf{Завдання 3}
\end{center}

\begin{verbatim}
print(not False>=4 or 3>7.0 and 3.0<=9)
\end{verbatim}



\begin{center}
\textbf{Завдання 5}
\end{center}

\begin{verbatim}
print(8.0 == False < 4 <= 8)
\end{verbatim}



\begin{center}
\textbf{Завдання 8}
\end{center}

\begin{verbatim}
a, b, c = 6, 7, 8
def g(a, b, c=9):
    print(a, b, c, end=" ")

g(b=5, c=1, 5)
print(a, b, c)
\end{verbatim}



\begin{center}
\textbf{Завдання 9}
\end{center}

\begin{verbatim}
a,b,c=1,5,4
def f(b):
    a=4
    b-=1
    c=3
    return a+b+c

a,b,c=6,3,8
print(f(a),a,b,c)
\end{verbatim}



\begin{center}
\textbf{Завдання 10}
\end{center}

Написати функцiю f, що отримує число і повертає логічну ознаку того, що воно не належить вiдрiзку [2017; 2018]. Стиль запису умови також оцінюється.\par

\end{minipage}

\end {large}}
\newpage

{\Large{06.11.2025 Контрольна робота \No 1, варіант  \No \textbf { 30 }\par}}
{
\begin {large}
У завданнях 1-9 вказати, що буде виведено за виконання.

Повідомлення інтерпретатора про помилку позначати ERROR

\begin{minipage}[t]{7.5cm}

\begin{center}
\textbf{Завдання 1}
\end{center}

\begin{verbatim}
print(6/8%3*4)
\end{verbatim}



\begin{center}
\textbf{Завдання 2}
\end{center}

\begin{verbatim}
print(-2**False+False)
\end{verbatim}



\begin{center}
\textbf{Завдання 4}
\end{center}

\begin{verbatim}
print("False" > 9)
\end{verbatim}



\begin{center}
\textbf{Завдання 6}
\end{center}

\begin{verbatim}

def f(n):
    print("f", end="");
    return n>=-2

print(f(1) and f(3))
\end{verbatim}



\begin{center}
\textbf{Завдання 7}
\end{center}

\begin{verbatim}
def h(a):
    u=75
    if a: 
        u=9
    elif a>-4:
         return 4
    else:
         u=3
    return u

print(h(-5))
\end{verbatim}

\end{minipage}
\vline \hspace{6mm}
\begin{minipage}[t]{10.5cm}

\begin{center}
\textbf{Завдання 3}
\end{center}

\begin{verbatim}
print(not 5<7 and 9.0!=True or 8.0==4)
\end{verbatim}



\begin{center}
\textbf{Завдання 5}
\end{center}

\begin{verbatim}
print(4.0 != True > 9 is 2.0)
\end{verbatim}



\begin{center}
\textbf{Завдання 8}
\end{center}

\begin{verbatim}
a, b, c = 9, 7, 6
def f(a, b, c):
    print(a, b, c, end=" ")

f(0, 2, 4)
print(a, b, c)
\end{verbatim}



\begin{center}
\textbf{Завдання 9}
\end{center}

\begin{verbatim}
a,b,c=9,7,0
def h(a):
    a=5
    b+=2
    c=2
    return a+b+c

a,b,c=2,4,9
print(h(b),a,b,c)
\end{verbatim}



\begin{center}
\textbf{Завдання 10}
\end{center}

Написати функцiю f, що отримує 2 числа і повертає логічну ознаку того, що їх сума дiлиться на 3. Стиль запису умови також оцінюється.\par

\end{minipage}

\end {large}}
\newpage

{\Large{06.11.2025 Контрольна робота \No 1, варіант  \No \textbf { 31 }\par}}
{
\begin {large}
У завданнях 1-9 вказати, що буде виведено за виконання.

Повідомлення інтерпретатора про помилку позначати ERROR

\begin{minipage}[t]{7.5cm}

\begin{center}
\textbf{Завдання 1}
\end{center}

\begin{verbatim}
print(5//7//9*5)
\end{verbatim}



\begin{center}
\textbf{Завдання 2}
\end{center}

\begin{verbatim}
print(1**1.0*True)
\end{verbatim}



\begin{center}
\textbf{Завдання 4}
\end{center}

\begin{verbatim}
print(2 > "6.0j")
\end{verbatim}



\begin{center}
\textbf{Завдання 6}
\end{center}

\begin{verbatim}

def f(n):
    print("f", end="");
    return n<4

print(f(-4) and f(2))
\end{verbatim}



\begin{center}
\textbf{Завдання 7}
\end{center}

\begin{verbatim}
def h(a,b):
    c=72
    if a:
        c=0
    elif a>=-3:
         return 7
    else: 
        c=8
    return c

print(h(3,-8))
\end{verbatim}

\end{minipage}
\vline \hspace{6mm}
\begin{minipage}[t]{10.5cm}

\begin{center}
\textbf{Завдання 3}
\end{center}

\begin{verbatim}
print(6.0>=5.0 and not 6>8 or 5==False)
\end{verbatim}



\begin{center}
\textbf{Завдання 5}
\end{center}

\begin{verbatim}
print(8.0 >= 7 != 2 < 7.0)
\end{verbatim}



\begin{center}
\textbf{Завдання 8}
\end{center}

\begin{verbatim}
a, b, c = 8, 8, 9
def g(a, b=6, c=7):
    print(a, b, c, end=" ")

g(1, c=3, b=3)
print(a, b, c)
\end{verbatim}



\begin{center}
\textbf{Завдання 9}
\end{center}

\begin{verbatim}
a,b,c=0,3,1
def g(a):
    a-=3
    b=5
    c=4
    return a+b+c

a,b,c=2,5,8
print(g(b),a,b,c)
\end{verbatim}



\begin{center}
\textbf{Завдання 10}
\end{center}

Написати функцiю f, що отримує число і повертає логічну ознаку того, що воно не належить iнтервалу (2020; 2021). Стиль запису умови також оцінюється.\par

\end{minipage}

\end {large}}
\newpage

{\Large{06.11.2025 Контрольна робота \No 1, варіант  \No \textbf { 32 }\par}}
{
\begin {large}
У завданнях 1-9 вказати, що буде виведено за виконання.

Повідомлення інтерпретатора про помилку позначати ERROR

\begin{minipage}[t]{7.5cm}

\begin{center}
\textbf{Завдання 1}
\end{center}

\begin{verbatim}
print(8/3/6*9)
\end{verbatim}



\begin{center}
\textbf{Завдання 2}
\end{center}

\begin{verbatim}
print(9**0.5*-2)
\end{verbatim}



\begin{center}
\textbf{Завдання 4}
\end{center}

\begin{verbatim}
print(False >= 4.0)
\end{verbatim}



\begin{center}
\textbf{Завдання 6}
\end{center}

\begin{verbatim}

def f(n):
    print("f", end="");
    return n

print(f(-2) or f(-8))
\end{verbatim}



\begin{center}
\textbf{Завдання 7}
\end{center}

\begin{verbatim}
def f(c):
    x=29
    if c>=-1: 
        x=5
    if c!=0:
         return 7
    else:
         return 1
    return x

print(f(-7))
\end{verbatim}

\end{minipage}
\vline \hspace{6mm}
\begin{minipage}[t]{10.5cm}

\begin{center}
\textbf{Завдання 3}
\end{center}

\begin{verbatim}
print(not 2.0<3 or True!=4.0 and 9<=2)
\end{verbatim}



\begin{center}
\textbf{Завдання 5}
\end{center}

\begin{verbatim}
print(False == 6 is 3 >= 3.0)
\end{verbatim}



\begin{center}
\textbf{Завдання 8}
\end{center}

\begin{verbatim}
a, b, c = 6, 9, 8
def f(a, b=7, c):
    print(a, b, c, end=" ")

f(b=5, a=0, c=1)
print(a, b, c)
\end{verbatim}



\begin{center}
\textbf{Завдання 9}
\end{center}

\begin{verbatim}
a,b,c=6,7,4
def f(b):
    a*=1
    b=1
    c=3
    return a+b+c

a,b,c=5,8,2
print(f(a),a,b,c)
\end{verbatim}



\begin{center}
\textbf{Завдання 10}
\end{center}

Написати функцiю f, що отримує 3 числа і повертає логічну ознаку того, що вони попарно рiзнi. Стиль запису умови також оцінюється.\par

\end{minipage}

\end {large}}
\newpage

{\Large{06.11.2025 Контрольна робота \No 1, варіант  \No \textbf { 33 }\par}}
{
\begin {large}
У завданнях 1-9 вказати, що буде виведено за виконання.

Повідомлення інтерпретатора про помилку позначати ERROR

\begin{minipage}[t]{7.5cm}

\begin{center}
\textbf{Завдання 1}
\end{center}

\begin{verbatim}
print(9//9%7*3)
\end{verbatim}



\begin{center}
\textbf{Завдання 2}
\end{center}

\begin{verbatim}
print(4**-0.5*-1)
\end{verbatim}



\begin{center}
\textbf{Завдання 4}
\end{center}

\begin{verbatim}
print("True" <= "4")
\end{verbatim}



\begin{center}
\textbf{Завдання 6}
\end{center}

\begin{verbatim}

def f(n):
    print("f", end="");
    return n

print(f(4) and f(-1))
\end{verbatim}



\begin{center}
\textbf{Завдання 7}
\end{center}

\begin{verbatim}
def f(a,b):
    c=87
    if b<-4:
        return 9
    if b==-5:
         c=6
    else: 
        return 5
    return c

print(f(-6,9))
\end{verbatim}

\end{minipage}
\vline \hspace{6mm}
\begin{minipage}[t]{10.5cm}

\begin{center}
\textbf{Завдання 3}
\end{center}

\begin{verbatim}
print(not 7.0==4 and False>4.0 or 7<2)
\end{verbatim}



\begin{center}
\textbf{Завдання 5}
\end{center}

\begin{verbatim}
print(9 <= True > 4.0 <= 6)
\end{verbatim}



\begin{center}
\textbf{Завдання 8}
\end{center}

\begin{verbatim}
a, b, c = 9, 6, 7
def h(a, b=8, c=6):
    print(a, b, c, end=" ")

h(4, b=2)
print(a, b, c)
\end{verbatim}



\begin{center}
\textbf{Завдання 9}
\end{center}

\begin{verbatim}
a,b,c=9,7,1
def g(b):
    a=4
    b+=5
    c=2
    return a+b+c

a,b,c=6,0,3
print(g(a),a,b,c)
\end{verbatim}



\begin{center}
\textbf{Завдання 10}
\end{center}

Написати функцiю f, що отримує рядок i повертає логічну ознаку того, що вiн складається рівно з однієї великої латинської літери. Стиль запису умови також оцінюється.\par

\end{minipage}

\end {large}}
\newpage

{\Large{06.11.2025 Контрольна робота \No 1, варіант  \No \textbf { 34 }\par}}
{
\begin {large}
У завданнях 1-9 вказати, що буде виведено за виконання.

Повідомлення інтерпретатора про помилку позначати ERROR

\begin{minipage}[t]{7.5cm}

\begin{center}
\textbf{Завдання 1}
\end{center}

\begin{verbatim}
print(4/4/5*6)
\end{verbatim}



\begin{center}
\textbf{Завдання 2}
\end{center}

\begin{verbatim}
print(-3**True-False)
\end{verbatim}



\begin{center}
\textbf{Завдання 4}
\end{center}

\begin{verbatim}
print(2.0 < "1.0")
\end{verbatim}



\begin{center}
\textbf{Завдання 6}
\end{center}

\begin{verbatim}

def f(n):
    print("f", end="");
    return n>0

print(f(-5) or f(0))
\end{verbatim}



\begin{center}
\textbf{Завдання 7}
\end{center}

\begin{verbatim}
def g(b):
    y=43
    if b<=-2: 
        y=6
    elif b==3:
         y=2
    else:
         return 5
    return y

print(g(1))
\end{verbatim}

\end{minipage}
\vline \hspace{6mm}
\begin{minipage}[t]{10.5cm}

\begin{center}
\textbf{Завдання 3}
\end{center}

\begin{verbatim}
print(not 5<=9 or 9.0!=True and 2.0>=8)
\end{verbatim}



\begin{center}
\textbf{Завдання 5}
\end{center}

\begin{verbatim}
print(2.0 < 2 != True > 9.0)
\end{verbatim}



\begin{center}
\textbf{Завдання 8}
\end{center}

\begin{verbatim}
a, b, c = 9, 7, 8
def f(a, b, c=9):
    print(a, b, c, end=" ")

f(b=5, c=3, 2)
print(a, b, c)
\end{verbatim}



\begin{center}
\textbf{Завдання 9}
\end{center}

\begin{verbatim}
a,b,c=9,6,0
def h(a):
    a=2
    b*=1
    c=4
    return a+b+c

a,b,c=1,5,8
print(h(b),a,b,c)
\end{verbatim}



\begin{center}
\textbf{Завдання 10}
\end{center}

Написати функцiю f, що отримує 2 числа і повертає логічну ознаку того, що вони дають однакову остачу при дiленнi на 5. Стиль запису умови також оцінюється.\par

\end{minipage}

\end {large}}
\newpage

{\Large{06.11.2025 Контрольна робота \No 1, варіант  \No \textbf { 35 }\par}}
{
\begin {large}
У завданнях 1-9 вказати, що буде виведено за виконання.

Повідомлення інтерпретатора про помилку позначати ERROR

\begin{minipage}[t]{7.5cm}

\begin{center}
\textbf{Завдання 1}
\end{center}

\begin{verbatim}
print(3//5//4*7)
\end{verbatim}



\begin{center}
\textbf{Завдання 2}
\end{center}

\begin{verbatim}
print(1**1.0*1)
\end{verbatim}



\begin{center}
\textbf{Завдання 4}
\end{center}

\begin{verbatim}
print(6 != "7j")
\end{verbatim}



\begin{center}
\textbf{Завдання 6}
\end{center}

\begin{verbatim}

def f(n):
    print("f", end="");
    return n

print(f(-9) and f(6))
\end{verbatim}



\begin{center}
\textbf{Завдання 7}
\end{center}

\begin{verbatim}
def g(a,b):
    c=10
    if a!=5:
        return 2
    if a<=4:
         c=3
    else: 
        c=6
    return c

print(g(-8,-9))
\end{verbatim}

\end{minipage}
\vline \hspace{6mm}
\begin{minipage}[t]{10.5cm}

\begin{center}
\textbf{Завдання 3}
\end{center}

\begin{verbatim}
print(3.0<=6.0 or not 3>False and 6!=4)
\end{verbatim}



\begin{center}
\textbf{Завдання 5}
\end{center}

\begin{verbatim}
print(7 >= 8 is 5 == 5.0)
\end{verbatim}



\begin{center}
\textbf{Завдання 8}
\end{center}

\begin{verbatim}
a, b, c = 6, 7, 8
def h(a, b, c):
    print(a, b, c, end=" ")

h(4, 1)
print(a, b, c)
\end{verbatim}



\begin{center}
\textbf{Завдання 9}
\end{center}

\begin{verbatim}
a,b,c=2,4,3
def f(a):
    a=3
    b-=5
    c=3
    return a+b+c

a,b,c=7,3,4
print(f(a),a,b,c)
\end{verbatim}



\begin{center}
\textbf{Завдання 10}
\end{center}

Написати функцiю f, що отримує число і повертає логічну ознаку того, що воно не належить iнтервалу (2018; 2019). Стиль запису умови також оцінюється.\par

\end{minipage}

\end {large}}
\newpage

{\Large{06.11.2025 Контрольна робота \No 1, варіант  \No \textbf { 36 }\par}}
{
\begin {large}
У завданнях 1-9 вказати, що буде виведено за виконання.

Повідомлення інтерпретатора про помилку позначати ERROR

\begin{minipage}[t]{7.5cm}

\begin{center}
\textbf{Завдання 1}
\end{center}

\begin{verbatim}
print(9/4*3*3)
\end{verbatim}



\begin{center}
\textbf{Завдання 2}
\end{center}

\begin{verbatim}
print(2**-2+2)
\end{verbatim}



\begin{center}
\textbf{Завдання 4}
\end{center}

\begin{verbatim}
print("False" == True)
\end{verbatim}



\begin{center}
\textbf{Завдання 6}
\end{center}

\begin{verbatim}

def f(n):
    print("f", end="");
    return n<=1

print(f(-7) or f(9))
\end{verbatim}



\begin{center}
\textbf{Завдання 7}
\end{center}

\begin{verbatim}
def h(a,b):
    c=82
    if a:
        return 9
    elif b>3:
         return 5
    else: 
        c=3
    return c

print(h(-1,8))
\end{verbatim}

\end{minipage}
\vline \hspace{6mm}
\begin{minipage}[t]{10.5cm}

\begin{center}
\textbf{Завдання 3}
\end{center}

\begin{verbatim}
print(not 7==5.0 and True<8.0 or 8>=3)
\end{verbatim}



\begin{center}
\textbf{Завдання 5}
\end{center}

\begin{verbatim}
print(False == 3 >= 4 > 6.0)
\end{verbatim}



\begin{center}
\textbf{Завдання 8}
\end{center}

\begin{verbatim}
a, b, c = 8, 9, 7
def g(a, b=6, c):
    print(a, b, c, end=" ")

g(0, 0, a=3)
print(a, b, c)
\end{verbatim}



\begin{center}
\textbf{Завдання 9}
\end{center}

\begin{verbatim}
a,b,c=0,8,9
def h(b):
    a=4
    b+=2
    c=5
    return a+b+c

a,b,c=7,6,5
print(h(a),a,b,c)
\end{verbatim}



\begin{center}
\textbf{Завдання 10}
\end{center}

Написати функцiю f, що отримує рядок i повертає логічну ознаку того, що вiн складається рівно з однієї літери, яка є десятковою цифрою. Стиль запису умови також оцінюється.\par

\end{minipage}

\end {large}}
\newpage

{\Large{06.11.2025 Контрольна робота \No 1, варіант  \No \textbf { 37 }\par}}
{
\begin {large}
У завданнях 1-9 вказати, що буде виведено за виконання.

Повідомлення інтерпретатора про помилку позначати ERROR

\begin{minipage}[t]{7.5cm}

\begin{center}
\textbf{Завдання 1}
\end{center}

\begin{verbatim}
print(7//9*4*6)
\end{verbatim}



\begin{center}
\textbf{Завдання 2}
\end{center}

\begin{verbatim}
print(4**0.5*-1)
\end{verbatim}



\begin{center}
\textbf{Завдання 4}
\end{center}

\begin{verbatim}
print("True" == 1)
\end{verbatim}



\begin{center}
\textbf{Завдання 6}
\end{center}

\begin{verbatim}

def f(n):
    print("f", end="");
    return n

print(f(7) or f(8))
\end{verbatim}



\begin{center}
\textbf{Завдання 7}
\end{center}

\begin{verbatim}
def h(d):
    v=37
    if d: 
        v=0
    elif d<=4:
         return 8
    else:
         return 9
    return v

print(h(-4))
\end{verbatim}

\end{minipage}
\vline \hspace{6mm}
\begin{minipage}[t]{10.5cm}

\begin{center}
\textbf{Завдання 3}
\end{center}

\begin{verbatim}
print(6.0>=True or not 9==2 and 2.0<5)
\end{verbatim}



\begin{center}
\textbf{Завдання 5}
\end{center}

\begin{verbatim}
print(3 != 6.0 < 6 is True)
\end{verbatim}



\begin{center}
\textbf{Завдання 8}
\end{center}

\begin{verbatim}
a, b, c = 8, 7, 9
def h(a, b, c):
    print(a, b, c, end=" ")

h(1, c=2, b=4)
print(a, b, c)
\end{verbatim}



\begin{center}
\textbf{Завдання 9}
\end{center}

\begin{verbatim}
a,b,c=0,3,1
def g(a):
    a=3
    b=4
    c=1
    return a+b+c

a,b,c=0,9,7
print(g(b),a,b,c)
\end{verbatim}



\begin{center}
\textbf{Завдання 10}
\end{center}

Написати функцiю f, що отримує рядок i повертає логічну ознаку того, що вiн складається рівно з однієї літери, яка є вісімковою цифрою. Стиль запису умови також оцінюється.\par

\end{minipage}

\end {large}}
\newpage

{\Large{06.11.2025 Контрольна робота \No 1, варіант  \No \textbf { 38 }\par}}
{
\begin {large}
У завданнях 1-9 вказати, що буде виведено за виконання.

Повідомлення інтерпретатора про помилку позначати ERROR

\begin{minipage}[t]{7.5cm}

\begin{center}
\textbf{Завдання 1}
\end{center}

\begin{verbatim}
print(8/5//9*4)
\end{verbatim}



\begin{center}
\textbf{Завдання 2}
\end{center}

\begin{verbatim}
print(1**-1.0*-2)
\end{verbatim}



\begin{center}
\textbf{Завдання 4}
\end{center}

\begin{verbatim}
print(False >= 3.0)
\end{verbatim}



\begin{center}
\textbf{Завдання 6}
\end{center}

\begin{verbatim}

def f(n):
    print("f", end="");
    return n==-1

print(f(-3) and f(5))
\end{verbatim}



\begin{center}
\textbf{Завдання 7}
\end{center}

\begin{verbatim}
def g(a,b):
    c=60
    if b!=1:
        c=7
    if b>-2:
         return 1
    else: 
        c=4
    return c

print(g(-3,-1))
\end{verbatim}

\end{minipage}
\vline \hspace{6mm}
\begin{minipage}[t]{10.5cm}

\begin{center}
\textbf{Завдання 3}
\end{center}

\begin{verbatim}
print(6>False and not 9.0<=7 or 8!=8.0)
\end{verbatim}



\begin{center}
\textbf{Завдання 5}
\end{center}

\begin{verbatim}
print(2.0 <= 8 > 4.0 <= 9)
\end{verbatim}



\begin{center}
\textbf{Завдання 8}
\end{center}

\begin{verbatim}
a, b, c = 6, 7, 6
def g(a, b, c=9):
    print(a, b, c, end=" ")

g(c=5, b=3, c=2)
print(a, b, c)
\end{verbatim}



\begin{center}
\textbf{Завдання 9}
\end{center}

\begin{verbatim}
a,b,c=2,1,4
def f(a):
    a=1
    b*=4
    c=3
    return a+b+c

a,b,c=3,2,9
print(f(b),a,b,c)
\end{verbatim}



\begin{center}
\textbf{Завдання 10}
\end{center}

Написати функцiю f, що отримує число і повертає логічну ознаку того, що воно не належить iнтервалу (2018; 2019). Стиль запису умови також оцінюється.\par

\end{minipage}

\end {large}}
\newpage

{\Large{06.11.2025 Контрольна робота \No 1, варіант  \No \textbf { 39 }\par}}
{
\begin {large}
У завданнях 1-9 вказати, що буде виведено за виконання.

Повідомлення інтерпретатора про помилку позначати ERROR

\begin{minipage}[t]{7.5cm}

\begin{center}
\textbf{Завдання 1}
\end{center}

\begin{verbatim}
print(4/6/6*7)
\end{verbatim}



\begin{center}
\textbf{Завдання 2}
\end{center}

\begin{verbatim}
print(-2**-2+2)
\end{verbatim}



\begin{center}
\textbf{Завдання 4}
\end{center}

\begin{verbatim}
print("5" != "7.0j")
\end{verbatim}



\begin{center}
\textbf{Завдання 6}
\end{center}

\begin{verbatim}

def f(n):
    print("f", end="");
    return n

print(f(2) or f(7))
\end{verbatim}



\begin{center}
\textbf{Завдання 7}
\end{center}

\begin{verbatim}
def f(a):
    z=66
    if a>=2: 
        z=7
    if a<-3:
         z=3
    else:
         return 4
    return z

print(f(-2))
\end{verbatim}

\end{minipage}
\vline \hspace{6mm}
\begin{minipage}[t]{10.5cm}

\begin{center}
\textbf{Завдання 3}
\end{center}

\begin{verbatim}
print(not 2<False or 5<=3.0 and 6!=7.0)
\end{verbatim}



\begin{center}
\textbf{Завдання 5}
\end{center}

\begin{verbatim}
print(5.0 < 7 >= False == 2)
\end{verbatim}



\begin{center}
\textbf{Завдання 8}
\end{center}

\begin{verbatim}
a, b, c = 8, 8, 7
def f(a, b=6, c=9):
    print(a, b, c, end=" ")

f(a=1, 4, a=5)
print(a, b, c)
\end{verbatim}



\begin{center}
\textbf{Завдання 9}
\end{center}

\begin{verbatim}
a,b,c=6,8,7
def g(b):
    a=2
    b-=5
    c=1
    return a+b+c

a,b,c=0,1,5
print(g(b),a,b,c)
\end{verbatim}



\begin{center}
\textbf{Завдання 10}
\end{center}

Написати функцiю f, що отримує 2 числа і повертає логічну ознаку того, що вони рiвнi з точнiстю до 0.01. Стиль запису умови також оцінюється.\par

\end{minipage}

\end {large}}
\newpage

{\Large{06.11.2025 Контрольна робота \No 1, варіант  \No \textbf { 40 }\par}}
{
\begin {large}
У завданнях 1-9 вказати, що буде виведено за виконання.

Повідомлення інтерпретатора про помилку позначати ERROR

\begin{minipage}[t]{7.5cm}

\begin{center}
\textbf{Завдання 1}
\end{center}

\begin{verbatim}
print(6//7%8*5)
\end{verbatim}



\begin{center}
\textbf{Завдання 2}
\end{center}

\begin{verbatim}
print(3**True-1)
\end{verbatim}



\begin{center}
\textbf{Завдання 4}
\end{center}

\begin{verbatim}
print("8j" < "False")
\end{verbatim}



\begin{center}
\textbf{Завдання 6}
\end{center}

\begin{verbatim}

def f(n):
    print("f", end="");
    return n!=4

print(f(-9) and f(-8))
\end{verbatim}



\begin{center}
\textbf{Завдання 7}
\end{center}

\begin{verbatim}
def f(a,b):
    c=58
    if a:
        return 0
    elif a<0:
         return 2
    else: 
        c=8
    return c

print(f(2,7))
\end{verbatim}

\end{minipage}
\vline \hspace{6mm}
\begin{minipage}[t]{10.5cm}

\begin{center}
\textbf{Завдання 3}
\end{center}

\begin{verbatim}
print(not 5.0>True and 4>=4.0 or 9==3)
\end{verbatim}



\begin{center}
\textbf{Завдання 5}
\end{center}

\begin{verbatim}
print(5 != 3.0 is 8.0 <= 4)
\end{verbatim}



\begin{center}
\textbf{Завдання 8}
\end{center}

\begin{verbatim}
a, b, c = 7, 9, 8
def f(a, b=6, c):
    print(a, b, c, end=" ")

f(0, 2, 4)
print(a, b, c)
\end{verbatim}



\begin{center}
\textbf{Завдання 9}
\end{center}

\begin{verbatim}
a,b,c=4,2,3
def g(b):
    a=2
    b=3
    c=5
    return a+b+c

a,b,c=8,6,5
print(g(a),a,b,c)
\end{verbatim}



\begin{center}
\textbf{Завдання 10}
\end{center}

Написати функцiю f, що отримує число і повертає логічну ознаку того, що воно належить iнтервалу (2020; 2021). Стиль запису умови також оцінюється.\par

\end{minipage}

\end {large}}
\newpage

{\Large{06.11.2025 Контрольна робота \No 1, варіант  \No \textbf { 41 }\par}}
{
\begin {large}
У завданнях 1-9 вказати, що буде виведено за виконання.

Повідомлення інтерпретатора про помилку позначати ERROR

\begin{minipage}[t]{7.5cm}

\begin{center}
\textbf{Завдання 1}
\end{center}

\begin{verbatim}
print(5/3%5*9)
\end{verbatim}



\begin{center}
\textbf{Завдання 2}
\end{center}

\begin{verbatim}
print(1**-4+True)
\end{verbatim}



\begin{center}
\textbf{Завдання 4}
\end{center}

\begin{verbatim}
print("8.0" > True)
\end{verbatim}



\begin{center}
\textbf{Завдання 6}
\end{center}

\begin{verbatim}

def f(n):
    print("f", end="");
    return n

print(f(-7) and f(6))
\end{verbatim}



\begin{center}
\textbf{Завдання 7}
\end{center}

\begin{verbatim}
def g(c):
    u=17
    if c!=-3: 
        u=3
    if c==-2:
         return 5
    else:
         u=0
    return u

print(g(-4))
\end{verbatim}

\end{minipage}
\vline \hspace{6mm}
\begin{minipage}[t]{10.5cm}

\begin{center}
\textbf{Завдання 3}
\end{center}

\begin{verbatim}
print(5.0<6 and not 4<=9 and 8.0>=True)
\end{verbatim}



\begin{center}
\textbf{Завдання 5}
\end{center}

\begin{verbatim}
print(9 >= True > False is 2)
\end{verbatim}



\begin{center}
\textbf{Завдання 8}
\end{center}

\begin{verbatim}
a, b, c = 6, 8, 7
def g(a, b=9, c=7):
    print(a, b, c, end=" ")

g(5, 3, a=1)
print(a, b, c)
\end{verbatim}



\begin{center}
\textbf{Завдання 9}
\end{center}

\begin{verbatim}
a,b,c=5,8,0
def f(a):
    a=3
    b=2
    c=1
    return a+b+c

a,b,c=7,4,1
print(f(a),a,b,c)
\end{verbatim}



\begin{center}
\textbf{Завдання 10}
\end{center}

Написати функцiю f, що отримує 2 числа і повертає логічну ознаку того, що перше число бiльше за куб другого. Стиль запису умови також оцінюється.\par

\end{minipage}

\end {large}}
\newpage

{\Large{06.11.2025 Контрольна робота \No 1, варіант  \No \textbf { 42 }\par}}
{
\begin {large}
У завданнях 1-9 вказати, що буде виведено за виконання.

Повідомлення інтерпретатора про помилку позначати ERROR

\begin{minipage}[t]{7.5cm}

\begin{center}
\textbf{Завдання 1}
\end{center}

\begin{verbatim}
print(3//8//7*8)
\end{verbatim}



\begin{center}
\textbf{Завдання 2}
\end{center}

\begin{verbatim}
print(9**1.0*False)
\end{verbatim}



\begin{center}
\textbf{Завдання 4}
\end{center}

\begin{verbatim}
print(2.0 <= 9)
\end{verbatim}



\begin{center}
\textbf{Завдання 6}
\end{center}

\begin{verbatim}

def f(n):
    print("f", end="");
    return n<3

print(f(8) or f(-5))
\end{verbatim}



\begin{center}
\textbf{Завдання 7}
\end{center}

\begin{verbatim}
def h(a,b):
    c=44
    if a==0:
        c=0
    if b>=-5:
         return 7
    else: 
        c=9
    return c

print(h(-8,6))
\end{verbatim}

\end{minipage}
\vline \hspace{6mm}
\begin{minipage}[t]{10.5cm}

\begin{center}
\textbf{Завдання 3}
\end{center}

\begin{verbatim}
print(False>4.0 or not 2==3 or 7!=9.0)
\end{verbatim}



\begin{center}
\textbf{Завдання 5}
\end{center}

\begin{verbatim}
print(5 == 7.0 < 9.0 != 3)
\end{verbatim}



\begin{center}
\textbf{Завдання 8}
\end{center}

\begin{verbatim}
a, b, c = 9, 6, 8
def h(a, b=6, c):
    print(a, b, c, end=" ")

h(b=0, c=5, 1)
print(a, b, c)
\end{verbatim}



\begin{center}
\textbf{Завдання 9}
\end{center}

\begin{verbatim}
a,b,c=5,0,4
def g(b):
    a=3
    b+=2
    c=1
    return a+b+c

a,b,c=7,3,1
print(g(a),a,b,c)
\end{verbatim}



\begin{center}
\textbf{Завдання 10}
\end{center}

Написати функцiю f, що отримує 2 цілих числа i повертає логічну ознаку того, що їх дев'ятковий запис закінчується на одну ту саму цифру. Стиль запису умови також оцінюється.\par

\end{minipage}

\end {large}}
\newpage

{\Large{06.11.2025 Контрольна робота \No 1, варіант  \No \textbf { 43 }\par}}
{
\begin {large}
У завданнях 1-9 вказати, що буде виведено за виконання.

Повідомлення інтерпретатора про помилку позначати ERROR

\begin{minipage}[t]{7.5cm}

\begin{center}
\textbf{Завдання 1}
\end{center}

\begin{verbatim}
print(5//8*9*4)
\end{verbatim}



\begin{center}
\textbf{Завдання 2}
\end{center}

\begin{verbatim}
print(9**0.5*-2)
\end{verbatim}



\begin{center}
\textbf{Завдання 4}
\end{center}

\begin{verbatim}
print("True" <= False)
\end{verbatim}



\begin{center}
\textbf{Завдання 6}
\end{center}

\begin{verbatim}

def f(n):
    print("f", end="");
    return n

print(f(-4) or f(-2))
\end{verbatim}



\begin{center}
\textbf{Завдання 7}
\end{center}

\begin{verbatim}
def f(d):
    x=49
    if d>3: 
        return 2
    elif d!=0:
         x=6
    else:
         return 8
    return x

print(f(-9))
\end{verbatim}

\end{minipage}
\vline \hspace{6mm}
\begin{minipage}[t]{10.5cm}

\begin{center}
\textbf{Завдання 3}
\end{center}

\begin{verbatim}
print(not 3.0>=8 or 5>True or 6.0!=3)
\end{verbatim}



\begin{center}
\textbf{Завдання 5}
\end{center}

\begin{verbatim}
print(6 is 7 != False > 3.0)
\end{verbatim}



\begin{center}
\textbf{Завдання 8}
\end{center}

\begin{verbatim}
a, b, c = 7, 8, 9
def g(a, b, c=9):
    print(a, b, c, end=" ")

g(c=3, b=4, c=0)
print(a, b, c)
\end{verbatim}



\begin{center}
\textbf{Завдання 9}
\end{center}

\begin{verbatim}
a,b,c=9,6,2
def f(a):
    a*=4
    b=5
    c=3
    return a+b+c

a,b,c=8,8,4
print(f(b),a,b,c)
\end{verbatim}



\begin{center}
\textbf{Завдання 10}
\end{center}

Написати функцiю f, що отримує 2 числа і повертає логічну ознаку того, що перше бiльше. Стиль запису умови також оцінюється.\par

\end{minipage}

\end {large}}
\newpage

{\Large{06.11.2025 Контрольна робота \No 1, варіант  \No \textbf { 44 }\par}}
{
\begin {large}
У завданнях 1-9 вказати, що буде виведено за виконання.

Повідомлення інтерпретатора про помилку позначати ERROR

\begin{minipage}[t]{7.5cm}

\begin{center}
\textbf{Завдання 1}
\end{center}

\begin{verbatim}
print(6/4/7*9)
\end{verbatim}



\begin{center}
\textbf{Завдання 2}
\end{center}

\begin{verbatim}
print(4**-3-False)
\end{verbatim}



\begin{center}
\textbf{Завдання 4}
\end{center}

\begin{verbatim}
print("4j" >= 3.0)
\end{verbatim}



\begin{center}
\textbf{Завдання 6}
\end{center}

\begin{verbatim}

def f(n):
    print("f", end="");
    return n>=0

print(f(0) and f(-6))
\end{verbatim}



\begin{center}
\textbf{Завдання 7}
\end{center}

\begin{verbatim}
def g(b):
    y=35
    if b: 
        y=7
    if b>-1:
         return 4
    else:
         y=9
    return y

print(g(6))
\end{verbatim}

\end{minipage}
\vline \hspace{6mm}
\begin{minipage}[t]{10.5cm}

\begin{center}
\textbf{Завдання 3}
\end{center}

\begin{verbatim}
print(not 9<=2.0 and 7==7.0 and False<4)
\end{verbatim}



\begin{center}
\textbf{Завдання 5}
\end{center}

\begin{verbatim}
print(2.0 <= 4 == 8 >= True)
\end{verbatim}



\begin{center}
\textbf{Завдання 8}
\end{center}

\begin{verbatim}
a, b, c = 7, 8, 6
def h(a, b, c):
    print(a, b, c, end=" ")

h(2, 5, 4)
print(a, b, c)
\end{verbatim}



\begin{center}
\textbf{Завдання 9}
\end{center}

\begin{verbatim}
a,b,c=2,9,6
def h(b):
    a=5
    b=4
    c=2
    return a+b+c

a,b,c=3,6,5
print(h(b),a,b,c)
\end{verbatim}



\begin{center}
\textbf{Завдання 10}
\end{center}

Написати функцiю f, що отримує 2 числа і повертає логічну ознаку того, що вони дають однакову остачу при дiленнi на 5. Стиль запису умови також оцінюється.\par

\end{minipage}

\end {large}}
\newpage

{\Large{06.11.2025 Контрольна робота \No 1, варіант  \No \textbf { 45 }\par}}
{
\begin {large}
У завданнях 1-9 вказати, що буде виведено за виконання.

Повідомлення інтерпретатора про помилку позначати ERROR

\begin{minipage}[t]{7.5cm}

\begin{center}
\textbf{Завдання 1}
\end{center}

\begin{verbatim}
print(9/3//8*5)
\end{verbatim}



\begin{center}
\textbf{Завдання 2}
\end{center}

\begin{verbatim}
print(2**4+-1)
\end{verbatim}



\begin{center}
\textbf{Завдання 4}
\end{center}

\begin{verbatim}
print("4.0j" < 1)
\end{verbatim}



\begin{center}
\textbf{Завдання 6}
\end{center}

\begin{verbatim}

def f(n):
    print("f", end="");
    return n<=1

print(f(4) or f(1))
\end{verbatim}



\begin{center}
\textbf{Завдання 7}
\end{center}

\begin{verbatim}
def f(a,b):
    c=91
    if b<=-4:
        return 8
    elif a>2:
         c=5
    else: 
        return 6
    return c

print(f(7,-6))
\end{verbatim}

\end{minipage}
\vline \hspace{6mm}
\begin{minipage}[t]{10.5cm}

\begin{center}
\textbf{Завдання 3}
\end{center}

\begin{verbatim}
print(not 8.0<4.0 and False<=8 or 5>2)
\end{verbatim}



\begin{center}
\textbf{Завдання 5}
\end{center}

\begin{verbatim}
print(6.0 < 4.0 == 5 != 2)
\end{verbatim}



\begin{center}
\textbf{Завдання 8}
\end{center}

\begin{verbatim}
a, b, c = 6, 8, 9
def f(a, b, c):
    print(a, b, c, end=" ")

f(3, c=1, b=0)
print(a, b, c)
\end{verbatim}



\begin{center}
\textbf{Завдання 9}
\end{center}

\begin{verbatim}
a,b,c=0,1,3
def h(a):
    a=5
    b-=4
    c=2
    return a+b+c

a,b,c=9,2,5
print(h(a),a,b,c)
\end{verbatim}



\begin{center}
\textbf{Завдання 10}
\end{center}

Написати функцiю f, що отримує 2 цілих числа i повертає логічну ознаку того, що їх дванадцятковий запис закінчується на одну ту саму цифру. Стиль запису умови також оцінюється.\par

\end{minipage}

\end {large}}
\newpage

{\Large{06.11.2025 Контрольна робота \No 1, варіант  \No \textbf { 46 }\par}}
{
\begin {large}
У завданнях 1-9 вказати, що буде виведено за виконання.

Повідомлення інтерпретатора про помилку позначати ERROR

\begin{minipage}[t]{7.5cm}

\begin{center}
\textbf{Завдання 1}
\end{center}

\begin{verbatim}
print(3//5%6*6)
\end{verbatim}



\begin{center}
\textbf{Завдання 2}
\end{center}

\begin{verbatim}
print(1**-0.5*True)
\end{verbatim}



\begin{center}
\textbf{Завдання 4}
\end{center}

\begin{verbatim}
print("7.0" > False)
\end{verbatim}



\begin{center}
\textbf{Завдання 6}
\end{center}

\begin{verbatim}

def f(n):
    print("f", end="");
    return n

print(f(3) and f(-1))
\end{verbatim}



\begin{center}
\textbf{Завдання 7}
\end{center}

\begin{verbatim}
def h(a):
    w=88
    if a<4: 
        return 1
    elif a>=-4:
         return 4
    else:
         w=1
    return w

print(h(8))
\end{verbatim}

\end{minipage}
\vline \hspace{6mm}
\begin{minipage}[t]{10.5cm}

\begin{center}
\textbf{Завдання 3}
\end{center}

\begin{verbatim}
print(9.0==6 or not 6!=6.0 and 5>=True)
\end{verbatim}



\begin{center}
\textbf{Завдання 5}
\end{center}

\begin{verbatim}
print(False > 9 is 7.0 <= 5.0)
\end{verbatim}



\begin{center}
\textbf{Завдання 8}
\end{center}

\begin{verbatim}
a, b, c = 7, 8, 6
def h(a, b=9, c=7):
    print(a, b, c, end=" ")

h(2, 5)
print(a, b, c)
\end{verbatim}



\begin{center}
\textbf{Завдання 9}
\end{center}

\begin{verbatim}
a,b,c=6,7,9
def h(b):
    a*=1
    b=1
    c=5
    return a+b+c

a,b,c=5,1,4
print(h(a),a,b,c)
\end{verbatim}



\begin{center}
\textbf{Завдання 10}
\end{center}

Написати функцiю f, що отримує число і повертає логічну ознаку того, що воно не належить iнтервалу (2019; 2021). Стиль запису умови також оцінюється.\par

\end{minipage}

\end {large}}
\newpage

{\Large{06.11.2025 Контрольна робота \No 1, варіант  \No \textbf { 47 }\par}}
{
\begin {large}
У завданнях 1-9 вказати, що буде виведено за виконання.

Повідомлення інтерпретатора про помилку позначати ERROR

\begin{minipage}[t]{7.5cm}

\begin{center}
\textbf{Завдання 1}
\end{center}

\begin{verbatim}
print(7//9*3*8)
\end{verbatim}



\begin{center}
\textbf{Завдання 2}
\end{center}

\begin{verbatim}
print(4**-1.0*1)
\end{verbatim}



\begin{center}
\textbf{Завдання 4}
\end{center}

\begin{verbatim}
print("7" == "True")
\end{verbatim}



\begin{center}
\textbf{Завдання 6}
\end{center}

\begin{verbatim}

def f(n):
    print("f", end="");
    return n

print(f(9) or f(5))
\end{verbatim}



\begin{center}
\textbf{Завдання 7}
\end{center}

\begin{verbatim}
def g(a,b):
    c=85
    if a:
        return 4
    if b==4:
         c=3
    else: 
        c=2
    return c

print(g(-6,5))
\end{verbatim}

\end{minipage}
\vline \hspace{6mm}
\begin{minipage}[t]{10.5cm}

\begin{center}
\textbf{Завдання 3}
\end{center}

\begin{verbatim}
print(2.0>4 or not 9==True and 7.0>=7)
\end{verbatim}



\begin{center}
\textbf{Завдання 5}
\end{center}

\begin{verbatim}
print(7 >= 8 < 9.0 <= 4)
\end{verbatim}



\begin{center}
\textbf{Завдання 8}
\end{center}

\begin{verbatim}
a, b, c = 8, 9, 6
def f(a, b=7, c):
    print(a, b, c, end=" ")

f(0, a=3)
print(a, b, c)
\end{verbatim}



\begin{center}
\textbf{Завдання 9}
\end{center}

\begin{verbatim}
a,b,c=8,3,7
def f(b):
    a+=2
    b=4
    c=3
    return a+b+c

a,b,c=0,2,6
print(f(b),a,b,c)
\end{verbatim}



\begin{center}
\textbf{Завдання 10}
\end{center}

Написати функцiю f, що отримує 2 цілих числа i повертає логічну ознаку того, що їх сімковий запис закінчується на одну ту саму цифру. Стиль запису умови також оцінюється.\par

\end{minipage}

\end {large}}
\newpage

{\Large{06.11.2025 Контрольна робота \No 1, варіант  \No \textbf { 48 }\par}}
{
\begin {large}
У завданнях 1-9 вказати, що буде виведено за виконання.

Повідомлення інтерпретатора про помилку позначати ERROR

\begin{minipage}[t]{7.5cm}

\begin{center}
\textbf{Завдання 1}
\end{center}

\begin{verbatim}
print(8/6/5*7)
\end{verbatim}



\begin{center}
\textbf{Завдання 2}
\end{center}

\begin{verbatim}
print(-1**2.0-2)
\end{verbatim}



\begin{center}
\textbf{Завдання 4}
\end{center}

\begin{verbatim}
print(8.0 != 5)
\end{verbatim}



\begin{center}
\textbf{Завдання 6}
\end{center}

\begin{verbatim}

def f(n):
    print("f", end="");
    return n>-2

print(f(-3) and f(6))
\end{verbatim}



\begin{center}
\textbf{Завдання 7}
\end{center}

\begin{verbatim}
def g(a):
    u=42
    if a<=-5: 
        return 6
    if a==1:
         u=2
    else:
         return 9
    return u

print(g(9))
\end{verbatim}

\end{minipage}
\vline \hspace{6mm}
\begin{minipage}[t]{10.5cm}

\begin{center}
\textbf{Завдання 3}
\end{center}

\begin{verbatim}
print(3.0<=3 and not 8!=False or 5.0<2)
\end{verbatim}



\begin{center}
\textbf{Завдання 5}
\end{center}

\begin{verbatim}
print(8.0 >= True < 6 > 3)
\end{verbatim}



\begin{center}
\textbf{Завдання 8}
\end{center}

\begin{verbatim}
a, b, c = 8, 6, 9
def g(a, b, c=7):
    print(a, b, c, end=" ")

g(a=4, 1, b=2)
print(a, b, c)
\end{verbatim}



\begin{center}
\textbf{Завдання 9}
\end{center}

\begin{verbatim}
a,b,c=9,1,0
def g(a):
    a-=2
    b=4
    c=5
    return a+b+c

a,b,c=2,3,6
print(g(b),a,b,c)
\end{verbatim}



\begin{center}
\textbf{Завдання 10}
\end{center}

Написати функцiю f, що отримує 3 числа і повертає логічну ознаку того, що вони попарно рiзнi. Стиль запису умови також оцінюється.\par

\end{minipage}

\end {large}}
\newpage

{\Large{06.11.2025 Контрольна робота \No 1, варіант  \No \textbf { 49 }\par}}
{
\begin {large}
У завданнях 1-9 вказати, що буде виведено за виконання.

Повідомлення інтерпретатора про помилку позначати ERROR

\begin{minipage}[t]{7.5cm}

\begin{center}
\textbf{Завдання 1}
\end{center}

\begin{verbatim}
print(4/7%4*3)
\end{verbatim}



\begin{center}
\textbf{Завдання 2}
\end{center}

\begin{verbatim}
print(1**1.0*False)
\end{verbatim}



\begin{center}
\textbf{Завдання 4}
\end{center}

\begin{verbatim}
print("5.0j" >= False)
\end{verbatim}



\begin{center}
\textbf{Завдання 6}
\end{center}

\begin{verbatim}

def f(n):
    print("f", end="");
    return n

print(f(7) or f(-3))
\end{verbatim}



\begin{center}
\textbf{Завдання 7}
\end{center}

\begin{verbatim}
def h(a,b):
    c=23
    if a>=-2:
        return 1
    elif b<=-1:
         c=0
    else: 
        return 3
    return c

print(h(-5,9))
\end{verbatim}

\end{minipage}
\vline \hspace{6mm}
\begin{minipage}[t]{10.5cm}

\begin{center}
\textbf{Завдання 3}
\end{center}

\begin{verbatim}
print(not False!=2 or 7==6.0 or 5.0<=5)
\end{verbatim}



\begin{center}
\textbf{Завдання 5}
\end{center}

\begin{verbatim}
print(7.0 == 7 is False != 2.0)
\end{verbatim}



\begin{center}
\textbf{Завдання 8}
\end{center}

\begin{verbatim}
a, b, c = 6, 9, 8
def h(a, b, c):
    print(a, b, c, end=" ")

h(4, 1, 3)
print(a, b, c)
\end{verbatim}



\begin{center}
\textbf{Завдання 9}
\end{center}

\begin{verbatim}
a,b,c=8,5,4
def f(a):
    a=3
    b+=1
    c=2
    return a+b+c

a,b,c=7,3,7
print(f(a),a,b,c)
\end{verbatim}



\begin{center}
\textbf{Завдання 10}
\end{center}

Написати функцiю f, що отримує 2 числа і повертає логічну ознаку того, що їх сума дiлиться на 3. Стиль запису умови також оцінюється.\par

\end{minipage}

\end {large}}
\newpage

{\Large{06.11.2025 Контрольна робота \No 1, варіант  \No \textbf { 50 }\par}}
{
\begin {large}
У завданнях 1-9 вказати, що буде виведено за виконання.

Повідомлення інтерпретатора про помилку позначати ERROR

\begin{minipage}[t]{7.5cm}

\begin{center}
\textbf{Завдання 1}
\end{center}

\begin{verbatim}
print(9//8*3*6)
\end{verbatim}



\begin{center}
\textbf{Завдання 2}
\end{center}

\begin{verbatim}
print(-3**False-True)
\end{verbatim}



\begin{center}
\textbf{Завдання 4}
\end{center}

\begin{verbatim}
print(3 != 9.0)
\end{verbatim}



\begin{center}
\textbf{Завдання 6}
\end{center}

\begin{verbatim}

def f(n):
    print("f", end="");
    return n<-3

print(f(-8) and f(-9))
\end{verbatim}



\begin{center}
\textbf{Завдання 7}
\end{center}

\begin{verbatim}
def g(a,b):
    c=35
    if b!=1:
        return 5
    elif a<5:
         c=9
    else: 
        c=7
    return c

print(g(5,-8))
\end{verbatim}

\end{minipage}
\vline \hspace{6mm}
\begin{minipage}[t]{10.5cm}

\begin{center}
\textbf{Завдання 3}
\end{center}

\begin{verbatim}
print(9>4.0 and not True>=8 and 2.0<4)
\end{verbatim}



\begin{center}
\textbf{Завдання 5}
\end{center}

\begin{verbatim}
print(4 != 3 == 6 > True)
\end{verbatim}



\begin{center}
\textbf{Завдання 8}
\end{center}

\begin{verbatim}
a, b, c = 7, 7, 9
def g(a, b, c=6):
    print(a, b, c, end=" ")

g(5, a=0)
print(a, b, c)
\end{verbatim}



\begin{center}
\textbf{Завдання 9}
\end{center}

\begin{verbatim}
a,b,c=9,8,2
def h(b):
    a=1
    b*=3
    c=4
    return a+b+c

a,b,c=1,4,6
print(h(a),a,b,c)
\end{verbatim}



\begin{center}
\textbf{Завдання 10}
\end{center}

Написати функцiю f, що отримує 2 числа і повертає логічну ознаку того, що хоча б одне з них не нуль. Стиль запису умови також оцінюється.\par

\end{minipage}

\end {large}}
\newpage

%end
\end{document}

